% \iffalse meta-comment
% !TEX program = XeLaTeX
%
% DO NOT REMOVE THE MARK MANUALLY (used in l3build) -> %^^A [L3BUILD_REPLACED_MARK]
%
%<*internal>
\iffalse
%</internal>
%<*internal>
\fi
\begingroup
  \def\nameofLaTeXe{LaTeX2e}
\expandafter\endgroup\ifx\fmtname\nameofLaTeXe\else
\csname fi\endcsname
%</internal>
%<*install>
\input docstrip
\keepsilent
\askforoverwritefalse
\preamble

    Copyright (C) 2013-2014 by Xu Cheng <xucheng@me.com>
                  2014-     by hust-latex <https://github.com/hust-latex>

    This work may be distributed and/or modified under the
    conditions of the LaTeX Project Public License, either
    version 1.3c of this license or (at your option) any later
    version. This version of this license is in
       https://www.latex-project.org/lppl/lppl-1-3c.txt
    and the latest version of this license is in
       https://www.latex-project.org/lppl.txt
    and version 1.3c or later is part of all distributions of LaTeX
    version 2008 or later.

    This work has the LPPL maintenance status "maintained".

    The Current Maintainer of this work is HUANG Yuxi <hustthesis@hyxi.dev>.

\endpreamble
\postamble

    This work consists of the files \jobname.dtx
              and the derived files \jobname.ins,
                                    \jobname.cls,
                                    \jobname-doc.sty,
                                    \jobname-m.def,
                                    \jobname-d.def,
                                    \jobname.cbx,
                                    \jobname.bbx.

\endpostamble

\generate
  {
%</install>
%<*internal>
    \usedir{source/latex/\jobname}
    \file{\jobname.ins}{\from{\jobname.dtx}{install}}
%</internal>
%<*install>
    \usedir{tex/latex/\jobname}
    \file{\jobname.cls}{\from{\jobname.dtx}{class}}
    \file{\jobname-doc.sty}{\from{\jobname.dtx}{doc-sty}}
    \file{\jobname-m.def}{\from{\jobname.dtx}{def-m}}
    \file{\jobname-d.def}{\from{\jobname.dtx}{def-d}}
    \file{\jobname.cbx}{\from{\jobname.dtx}{cbx}}
    \file{\jobname.bbx}{\from{\jobname.dtx}{bbx}}
  }
\endbatchfile
%</install>
%<*internal>
\fi
%</internal>

%<class>\NeedsTeXFormat{LaTeX2e}[2024-11-01]
%<class>\ProvidesExplClass{hustthesis}
%<doc-sty>\ProvidesExplClass{hustthesis-doc}
%<def-m>\ProvidesExplFile{hustthesis-m.def}
%<def-d>\ProvidesExplFile{hustthesis-d.def}
%<cbx>\ProvidesFile{hustthesis.cbx}[Citation style for the hustthesis class]
%<bbx>\ProvidesFile{hustthesis.bbx}[Bibliography style for the hustthesis class]
%<class|doc-sty|def-m|def-d>  {[LATEST]}{[DEV]} %^^A [L3BUILD_REPLACED_MARK]
%<class>  {Thesis Template for Huazhong University of Science and Technology}
%<doc-sty>  {Documentation for the hustthesis class}
%<def-m>  {Definition file for the hustthesis class in master thesis}
%<def-d>  {Definition file for the hustthesis class in doctoral thesis}
%<*driver>
\documentclass[]{\jobname}
\usepackage{useclass}
% Protect \maketitile from the redefinition in \cls{article} in \cls{l3doc}.
\let\maketitleold=\maketitle
\useclass*[full]{l3doc}
\let\maketitle=\maketitleold
\usepackage{\jobname-doc}
\usepackage{mdframed}
\usepackage[backend=biber,style=hustthesis]{biblatex} % temporary
\begin{filecontents}[overwrite]{\jobname.bib}
@online{zeng2023fduthesis,
  title = {\cls{fduthesis}: 复旦大学论文模板},
  author = {曾, 祥东},
  year = {2023},
  month = {may},
  day = {27},
  url = {https://github.com/stone-zeng/fduthesis},
}
@online{sjtug2024SJTUTeX,
  title = {\cls{SJTU\TeX}: 上海交通大学文档类集},
  author = {SJTUG},
  year = {2024},
  month = {feb},
  day = {18},
  url = {https://github.com/sjtug/SJTUTeX},
}
@online{njulug2023njuthesis,
  title = {南京大学学位论文模板},
  author = {南京大学~Linux~用户组},
  year = {2023},
  month = {dec},
  day = {05},
  url = {https://github.com/nju-lug/NJUThesis},
}
\end{filecontents}
\addbibresource{\jobname.bib}

\providecommand{\name}{\jobname}
\providecommand{\clsname}{\cls{\name}}

\begin{document}
  \DocInput{\jobname.dtx}
  \PrintChanges
  \PrintIndex
\end{document}
%</driver>
% \fi
%
% \changes{v2.0}{2024-10-01}{使用\cls{useclass}在文档类手册中应用自身}
%
% \hustsetup {
%   info =
%     {
%       title = { \clsname{}：华中科技大学研究生\\学位论文\LaTeX{}文档类使用说明 },
%       title* = { \clsname{}: \LaTeX{} Thesis Template for Huazhong University of Science and Technology },
%       degree = { academic },
%       degree* = { Engineering },
%       author = { 姓名 },
%       author* = { Author },
%       student-id = { D201234567 },
%       major = { 专业 },
%       major* = { Major },
%       supervisor = { 指导老师\quad{} 教授 },
%       supervisor* = { Prof.~ Supervisor },
%     },
% }
% \title{The \clsname.cls Class}
% \author
%   {
%     Huang Yuxi^^A
%     \thanks{
%       \textit{Huang} is his surname.
%     }
%     <^^A
%     \href
%       {mailto:hustthesis@hyxi.dev}
%       {hustthesis@hyxi.dev}^^A
%     >
%   }
% \date{[LATEST] \quad Version [DEV]} ^^A [L3BUILD_REPLACED_MARK]
%
% \maketitle
%
% \tableofcontents
%
% \begin{documentation}
%
% \chapter{引言}
%
% 提前致谢 \cls{fduthesis} 文档类\cite{zeng2023fduthesis}， \cls{SJTU\TeX} 文档类集\cite{sjtug2024SJTUTeX}，和 \cls{njuthesis} 文档类\cite{njulug2023njuthesis}，本文档类的许多设计和实现都是参考了上述代码。
%
% \printbibliography
%
% \end{documentation}
%
% \begin{implementation}
%
% \newgeometry{
%   marginpar = 2.8 cm,
% }
%
% \chapter{\clsname{} 代码实现}
% 本文档类使用 \LaTeX3 语法编写，依赖 \pkg{expl3} 环境。
%
%    \begin{macrocode}
%<@@=hust>
%<*class>
%    \end{macrocode}
%
% \section{文档类预处理}
%
% 本文档类隐式依赖已整合进 \LaTeX--\pkg{base} 的 \pkg{ltkeys} (原 \pkg{l3keys2e}) 和 \pkg{lttemplates}（原 \pkg{xtemplate}）。
%
%    \begin{macrocode}
 \RequirePackage { }
%    \end{macrocode}
%
% \subsection{变量声明}
%
% \begin{variable}{
%   \l_@@_tmpa_box,
%   \l_@@_tmpa_clist,
%   \l_@@_tmpa_dim,
%   \l_@@_tmpa_skip,
%   \l_@@_tmpa_tl, \l_@@_tmpb_tl,
% }
% 临时变量。
%    \begin{macrocode}
\box_new:N \l_@@_tmpa_box
\clist_new:N \l_@@_tmpa_clist
\dim_new:N \l_@@_tmpa_dim
\skip_new:N \l_@@_tmpa_skip
\tl_new:N \l_@@_tmpa_tl
\tl_new:N \l_@@_tmpb_tl
%    \end{macrocode}
% \end{variable}
%
% \begin{variable}{\g_@@_degree_level_tl}
% 论文类型。
%    \begin{macrocode}
\tl_new:N \g_@@_degree_level_tl
%    \end{macrocode}
% \end{variable}
%
% \begin{variable}{\g_@@_to_ctex_class_clist, \g_@@_to_hyperref_clist}
% 分别保存由 \clsname{} 传入 \cls{ctexbook} 文档类和 \pkg{hyperref} 宏包的选项列表。
%    \begin{macrocode}
\clist_new:N \g_@@_to_ctex_class_clist
\clist_new:N \g_@@_to_hyperref_clist
%    \end{macrocode}
% \end{variable}
%
% \begin{variable}{\g_@@_draft_bool}
% 草稿模式。
%    \begin{macrocode}
\bool_new:N \g_@@_draft_bool
%    \end{macrocode}
% \end{variable}
%
% \begin{variable}{\g_@@_anonymous_bool}
% 盲审模式。
%    \begin{macrocode}
\bool_new:N \g_@@_anonymous_bool
%    \end{macrocode}
% \end{variable}
%
% \begin{variable}{
%   \g_@@_font_set_tl,
%   \g_@@_font_latin_tl,
%   \g_@@_font_cjk_tl,
%   \g_@@_font_math_tl}
% 存储所使用字体名称的全局变量。
%    \begin{macrocode}
\tl_new:N \g_@@_font_set_tl
\tl_new:N \g_@@_font_latin_tl
\tl_new:N \g_@@_font_cjk_tl
\tl_new:N \g_@@_font_math_tl
%    \end{macrocode}
% \end{variable}
%
% \begin{variable}{\g_@@_font_dir_tl}
% 存储字体路径的全局变量。
%    \begin{macrocode}
\tl_new:N \g_@@_font_dir_tl
%    \end{macrocode}
% \end{variable}
%
% \begin{variable}{\g_@@_font_dir_bool}
% 是否使用独立的字体文件。
%    \begin{macrocode}
\bool_new:N \g_@@_font_dir_bool
%    \end{macrocode}
% \end{variable}
%
% \begin{variable}{\g_@@_config_tl}
% 保存配置文件名称。默认为空。
%    \begin{macrocode}
\tl_new:N \g_@@_config_tl
%    \end{macrocode}
% \end{variable}
%
% \subsection{内部函数}
%
%
% \subsubsection{辅助函数}
% \begin{macro}{\@@_patch_cmd:Nnn, \@@_appto_cmd:Nn}
% 补丁工具，来自 \pkg{ctexpatch} 宏包。
%    \begin{macrocode}
\cs_new_protected:Npn \@@_patch_cmd:Nnn #1#2#3
  {
    \ctex_patch_cmd_once:NnnnTF #1 { } {#2} {#3}
      { } { \ctex_patch_failure:N #1 }
  }
\cs_new_protected:Npn \@@_appto_cmd:Nn #1#2
  {
    \ctex_appto_cmd:NnnTF #1 { } {#2}
      { } { \ctex_patch_failure:N #1 }
  }
%    \end{macrocode}
% \end{macro}
%
% \begin{macro}{\@@_define_fn_style:nn,
%   \@@_define_punct:nn,
%   \@@_define_name:nn,
%   \@@_define_name:nnn}
% 用来定义脚注样式、标点、默认名称的辅助函数。
%    \begin{macrocode}
\cs_new_protected:Npn \@@_define_fn_style:nn #1#2
  { \tl_const:cn { c_@@_fn_style_ #1 _tl } {#2} }
\cs_new_protected:Npn \@@_define_punct:nn #1#2
  { \tl_const:cn { c_@@_ #1 _tl } {#2} }
\cs_new_protected:Npn \@@_define_name:nn #1#2
  { \tl_const:cn { c_@@_name_ #1 _tl } {#2} }
\cs_new_protected:Npn \@@_define_name:nnn #1#2#3
  {
    \tl_const:cn { c_@@_name_ #1 _en_tl } {#2}
    \tl_const:cn { c_@@_name_ #1 _zh_tl } {#3}
  }
%    \end{macrocode}
% \end{macro}
%
%
% \begin{macro}{\@@_msg_new:nn,
%   \@@_error:n, \@@_error:nn, \@@_error:nx, \@@_error:nnn, \@@_error:nnnn,
%   \@@_warning:n, \@@_warning:nn, \@@_warning:nxx,
%   \@@_info:nx
%}
% 各种信息函数的缩略形式。
%    \begin{macrocode}
\cs_new:Npn \@@_msg_new:nn { \msg_new:nnn { hustthesis } }
\cs_new:Npn \@@_error:n { \msg_error:nn { hustthesis } }
\cs_new:Npn \@@_error:nn { \msg_error:nnn { hustthesis } }
\cs_new:Npn \@@_error:nx { \msg_error:nnx { hustthesis } }
\cs_new:Npn \@@_error:nnn { \msg_error:nnnn { hustthesis } }
\cs_new:Npn \@@_error:nnnn { \msg_error:nnnnn { hustthesis } }
\cs_new:Npn \@@_warning:n { \msg_warning:nn { hustthesis } }
\cs_new:Npn \@@_warning:nn { \msg_warning:nnn { hustthesis } }
\cs_new:Npn \@@_warning:nxx { \msg_warning:nnxx { hustthesis } }
\cs_new:Npn \@@_info:nx { \msg_info:nnx { hustthesis } }
%    \end{macrocode}
% \end{macro}
%
% \subsubsection{\LaTeXe{} 封装}
%
% \begin{macro}[deprecated]{\@@_vspace:N, \@@_vspace:n, \@@_vspace:c}
% 类似 \LaTeXe{} 中的 \tn{vspace*}。
%    \begin{macrocode}
\cs_new_protected:Npn \@@_vspace:N #1
  {
    \dim_set_eq:NN \l_@@_tmpa_dim \tex_prevdepth:D
    \tex_hrule:D height \c_zero_dim
    \nobreak
    \skip_vertical:N #1
    \skip_vertical:N \c_zero_skip
    \dim_set_eq:NN \tex_prevdepth:D \l_@@_tmpa_dim
  }
\cs_new_protected:Npn \@@_vspace:n #1
  {
    \skip_set:Nn \l_@@_tmpa_skip {#1}
    \@@_vspace:N \l_@@_tmpa_skip
  }
\cs_generate_variant:Nn \@@_vspace:N { c }
%    \end{macrocode}
% \end{macro}
%
% \begin{macro}[deprecated]{\@@_hrule:nn}
% 不改变位置的 \tn{hrule}。
% \begin{arguments}
%    \item 线段粗细（高度）
%    \item 线段长度（宽度）
% \end{arguments}
%    \begin{macrocode}
\cs_new_protected:Npn \@@_hrule:nn #1#2
  {
    \rule [ -0.5 ex ] { #2 } { #1 }
    \skip_horizontal:n { -#2 }
  }
%    \end{macrocode}
% \end{macro}
%
% \begin{macro}{\@@_ul_to_wd:nnn}
  % \begin{arguments}
  %    \item 下划线粗细
  %    \item 下划线长度
  %    \item 内容
  % \end{arguments}
%    \begin{macrocode}
\dim_new:N \l_@@_ul_dim
\cs_new_protected:Npn \@@_ul_to_wd:nnn #1#2#3
  {
    \hbox_overlap_right:n { \rule [ - 0.5 ex - #1 / 2 ] { #2 } { #1 } }
    \hbox_to_wd:nn { #2 } { \tex_hfil:D #3 \tex_hfil:D }
  }
%    \end{macrocode}
% \end{macro}
%
% \begin{macro}{\@@_line_spread:N, \@@_line_spread:n}
% 设置行距。
% \begin{arguments}
%   \item 行距倍数， |fp| 变量
% \end{arguments}
%    \begin{macrocode}
\cs_new:Npn \@@_line_spread:N #1
  {
    \tl_set:Nn \baselinestretch { \fp_use:N #1 }
    \selectfont
  }
\cs_new:Npn \@@_line_spread:n #1
  {
    \tl_set:Nn \baselinestretch { #1 }
    \selectfont
  }
%    \end{macrocode}
% \end{macro}
%
% \begin{macro}[deprecated]{\@@_arabic:n}
% 等价于 \LaTeXe{} 中的 \tn{arabic}。
%    \begin{macrocode}
\cs_new:Npn \@@_arabic:n #1
  { \int_to_arabic:v { c@ #1 } }
%    \end{macrocode}
% \end{macro}
%
% \subsection{选项处理}
%
% 定义 | hust / option | 键值类，该键值类仅可在文档类选项中使用。
%    \begin{macrocode}
\keys_define:nn { hust / option }
  {
%    \end{macrocode}
%
% \begin{macro}[module=option/]{level}
% 论文类型。目前分为硕士学位论文~\opt{master} 和博士学位论文~\opt{doctor}，默认为博士学位论文。
%    \begin{macrocode}
    level .choices:nn =
      { master, doctor }
      { \tl_gset_eq:NN \g_@@_degree_level_tl \l_keys_choice_tl },
    level .value_required:n = true,
    level .initial:n = { doctor },
%    \end{macrocode}
% \end{macro}
%
% \begin{macro}[module=option/]{draft}
% 是否启用草稿模式。默认为 \opt{false}。
%    \begin{macrocode}
    draft .choice:,
    draft / true .code:n =
      {
        \bool_gset_true:N \g_@@_draft_bool
        \clist_gput_right:Nn \g_@@_to_ctex_class_clist { draft }
      },
    draft / false .code:n =
      {
        \bool_gset_false:N \g_@@_draft_bool
      },
    draft .default:n = { true },
    draft .initial:n = { false },
%    \end{macrocode}
% \end{macro}
%
% \begin{macro}[module=option/]{anonymous}
% 是否启用盲审模式。默认为 \opt{false}，即显示作者信息。\key[option]{anonymous}
%    \begin{macrocode}
    anonymous .choice:,
    anonymous / true .code:n =
      {
        \bool_gset_true:N \g_@@_anonymous_bool
      },
    anonymous / false .code:n =
      {
        \bool_gset_false:N \g_@@_anonymous_bool
      },
    anonymous .default:n = { true },
    anonymous .initial:n = { false },
%    \end{macrocode}
% \end{macro}
%
% \begin{macro}[module=option/]{latin-font, cjk-font}
% 中英文字体选项。\opt{fandol} 和 \opt{gyre} 是等价的。
%    \begin{macrocode}
    latin-font .choices:nn =
      { fandol, gyre, mac, macoffice, win, none, }
      { \tl_gset:Nn \g_@@_font_latin_tl {#1} },
    cjk-font .choices:nn =
      { fandol, founder, mac, macoffice, noto, source, win, none, }
      { \tl_gset:Nn \g_@@_font_cjk_tl {#1} },
%    \end{macrocode}
% \end{macro}
%
% \begin{macro}[module=option/]{fontset}
% 中英文字体选项。\opt{fontset} 这个名称和 \pkg{ctex} 中的是一致的。
%    \begin{macrocode}
    fontset .choices:nn =
      { fandol, mac, macoffice, win, none, }
      { \keys_set:nn { hust / options } { latin-font = #1, cjk-font = #1 } },
%    \end{macrocode}
% \end{macro}
%
% \begin{macro}[module=option/]{math-font}
% 数学字体选项。
% 由 \pkg{unicode-math} 指定 \XeTeX 和 \LuaTeX 下使用的数学字体。
%    \begin{macrocode}
    math-font .choices:nn =
      {
        asana, cambria, fira, garamond, lm, libertinus, newcm,
        stix, bonum, dejavu, pagella, schola, termes, xits, none,
      }
      { \tl_gset_eq:NN \g_@@_font_math_tl \l_keys_choice_tl },
    math-font .initial:n = xits,
%    \end{macrocode}
% \end{macro}
%
% \begin{macro}[module=option/]{font-dir}
% 独立字体文件的路径。
%    \begin{macrocode}
    font-dir .code:n =
      {
        \bool_gset_true:N \g_@@_font_dir_bool
        \tl_gset_eq:NN \g_@@_font_dir_tl \l_keys_value_tl
      },
%    \end{macrocode}
% \end{macro}
%
% \begin{macro}[module=option/]{config}
% 自定义配置文件名。
%    \begin{macrocode}
    config .tl_gset:N = \g_@@_config_tl,
%    \end{macrocode}
% \end{macro}
%
% 将未知选项传给 \cls{ctexbook} 文档类。
%    \begin{macrocode}
    unknown .code:n =
      { \clist_gput_right:No \g_@@_to_ctex_class_clist { \CurrentOption } }
  }
%    \end{macrocode}
%
% \begin{macro}[module=option/, int]{option/}
% 将文档类选项传给 |hust / option|。
%    \begin{macrocode}
\ProcessKeyOptions [ hust / option ]
%    \end{macrocode}
% \end{macro}
%
% \subsection{常量载入}
%
% 载入常量定义文件。
%    \begin{macrocode}
\tl_if_empty:NTF \g_@@_config_tl
  {
    \tl_if_eq:NnT \g_@@_degree_level_tl { master }
      { \file_input:n { hustthesis-m.def } }
    \tl_if_eq:NnT \g_@@_degree_level_tl { doctor }
      { \file_input:n { hustthesis-d.def } }
  }
  {
    \@@_msg_new:nn { load-config-file }
      { You~ are~ loading~ config~ file~ "#1". }
    \@@_info:n { load-config-file }
    \file_input:V { \g_@@_config_tl }
  }
%    \end{macrocode}
%
% \subsection{宏包载入}
%
% \sout{对不支持粗体斜体的字体进行兼容处理，仅限于 \XeTeX{} 引擎。}\\
% \sout{参考自 \url{https://ask.latexstudio.net/ask/question/17464.html}。}
%    \begin{macrocode}
%\sys_if_engine_xetex:T
%  {
%    \PassOptionsToPackage { AutoFakeBold=4.0 } { xeCJK }
%    \AddToHook { package / ctex / before }
%      {
%        \AddToHook { package / xeCJK / after }
%          {
%            \keys_define:nn { xeCJK / features }
%              {
%                BoldFont .code:n = ,
%                ItalicFont .code:n =,
%              }
%          }
%      }
%    \AddToHook { package / ctex / after }
%      {
%        \keys_define:nn { xeCJK / features }
%          {
%            BoldFont .tl_set:N = \l__xeCJK_font_name_bf_tl ,
%            ItalicFont .tl_set:N = \l__xeCJK_font_name_it_tl,
%          }
%      }
%  }
%    \end{macrocode}
%
% 将选项传入 \cls{ctexbook} 文档类。
%    \begin{macrocode}
\PassOptionsToClass
  {
    a4paper,
    UTF8,
    scheme = chinese,
    fontset = none,
    oneside,
    heading = true,
    zihao = -4,
    linespread = \c_@@_line_spread_fp,
    \g_@@_to_ctex_class_clist,
  }
  { ctexbook }
\LoadClass { ctexbook }
%    \end{macrocode}
%
% 其中 \pkg{utfsym} 宏包不依赖外部字体，使用 \TikZ{} 宏包绘制 |U+2610|:~☐，|U+2611|:~☑，|U+1F5F9|:~\usym{1F5F9} 等符号。
%    \begin{macrocode}
\RequirePackage
  {
    amsmath,
    geometry,
    fancyhdr,
    titletoc,
    graphicx,
    longtable,
    caption,
    utfsym,
  }
%    \end{macrocode}
%
%
% \section{全局样式}
%
% \subsection{页面布局}
% 利用 \pkg{geometry} 宏包设置纸张大小、页面边距以及页眉高度。
%    \begin{macrocode}
\geometry
  {
    top = 4.5 cm,
    bottom = 2.8 cm,
    left = 2.8 cm,
    right = 2.8 cm,
    headheight = 30 pt,
  }
%    \end{macrocode}
%
% \subsection{字体设置}
%
% ^^A 本小节初版实现主要来自 \pkg{njuthesis} 文档类。
%
% \pkg{fontspec} 包含于 \pkg{ctex} 宏集，无需另行载入。
%
% \subsubsection{操作系统检测}
%
% \begin{variable}{\c_@@_path_macoffice_tl}
% mac Office 字体路径。
%    \begin{macrocode}
\tl_const:Nn \c_@@_path_macoffice_tl
  { /Applications/Microsoft~ Word.app/Contents/Resources/DFonts/ }
%    \end{macrocode}
% \end{variable}
%
% 操作系统检测。相较于 \pkg{ctex} 提供的 \cs{ctex_detect_platform:}，
% 额外增加了对 macOS 上 MS Office 的检测。
%    \begin{macrocode}
\sys_if_platform_windows:TF
  { \tl_gset:Nn \g_@@_font_set_tl { win } }
  {
    \ctex_if_platform_macos:TF
      {
        \file_if_exist:nTF { \c_@@_path_macoffice_tl times.ttf }
          { \tl_gset:Nn \g_@@_font_set_tl { macoffice } }
          { \tl_gset:Nn \g_@@_font_set_tl { mac } }
      }
      { \tl_gset:Nn \g_@@_font_set_tl { fandol } }
  }
%    \end{macrocode}
% 判断用户是否自定义了中英文字体。如果其中任意一种未被定义，
% 则使用系统预装字体覆盖字体选项。
% Windows 或 macOS 外的系统被判断为 Linux，一律使用自由字体。
%    \begin{macrocode}
\tl_if_empty:NT \g_@@_font_latin_tl
  { \tl_gset_eq:NN \g_@@_font_latin_tl \g_@@_font_set_tl }
\tl_if_empty:NT \g_@@_font_cjk_tl
  { \tl_gset_eq:NN \g_@@_font_cjk_tl \g_@@_font_set_tl }
%    \end{macrocode}
%
%
% \subsubsection{定义英文字库}
%
% 接下来逐个定义所需要使用的字库。
%
% \begin{macro}{
%   \@@_loadfont_latin:n,
%   \@@_loadfont_latin_win:,
%   \@@_loadfont_latin_mac:,
% }
% Windows 与 macOS 西文字体的区别主要在于默认等宽字体。
%    \begin{macrocode}
\cs_new:Npn \@@_loadfont_latin:n #1
  {
    \__fontspec_main_setmainfont:nn { } { Times~New~Roman }
    \__fontspec_main_setsansfont:nn { } { Arial }
    \__fontspec_main_setmonofont:nn { Scale = MatchLowercase } {#1}
  }
\cs_new_protected:Npn \@@_loadfont_latin_win:
  { \@@_loadfont_latin:n { Courier~New } }
\cs_new_protected:Npn \@@_loadfont_latin_mac:
  {
    \@@_loadfont_latin:n { Menlo }
%    \end{macrocode}
% 检测 Times New Roman 是否具有小型大写字母（small caps）字型。
% 这是 macOS 预装的字体版本较旧导致的。
%    \begin{macrocode}
    \fontspec_if_small_caps:F
      {
        \
        \@@_warning:n { no-small-caps }
        \__fontspec_main_setmainfont:nn
          { \c_@@_name_gyrefeature_clist } { texgyretermes }
      }
  }
\@@_msg_new:nn { no-small-caps }
  {
    I~ am~ using~ TeX~ Gyre~ Termes~ as~ default~ Roman~ font.\\
    This~ is~ because~ the~ "Times~ New~ Roman"~ font~ in~ your~
    system~ does~ not~ embed~ glyphs~ for~ small~ capitals.~
    You~ can~ ignore~ this~ warning~ if~ you~ do~ not~ need~
    \string\textsc.~ For~ more~ information,~
    please~ refer~ to~ section~ 3.2.6~ of~ the~ documentation.
  }
%    \end{macrocode}
% \end{macro}
%
% \begin{macro}{\@@_loadfont_latin_macoffice:}
% mac Office 西文字体。
%    \begin{macrocode}
\cs_new_protected:Npn \@@_loadfont_latin_macoffice:
  {
    \__fontspec_main_setmainfont:nn
      { \c_@@_name_macofficefeature_clist } { times }
    \__fontspec_main_setsansfont:nn
      { \c_@@_name_macofficefeature_clist } { arial }
    \__fontspec_main_setmonofont:nn
      { Scale = MatchLowercase } { Menlo }
  }
%    \end{macrocode}
% \end{macro}
%
% \begin{variable}{\c_@@_name_macofficefeature_clist}
% 用于 \pkg{fontspec} 的 mac Office 字体特性列表。
%    \begin{macrocode}
\clist_const:Nn \c_@@_name_macofficefeature_clist
  {
    Path = \c_@@_path_macoffice_tl,
    Extension = .ttf,
    UprightFont = *,
    BoldFont = *bd,
    ItalicFont = *i,
    BoldItalicFont = *bi,
  }
%    \end{macrocode}
% \end{variable}
%
% \begin{macro}{\@@_loadfont_latin_gyre:}
% 开源的 \TeX Gyre 西文字体。
%    \begin{macrocode}
\cs_new_protected:Npn \@@_loadfont_latin_gyre:
  {
    \__fontspec_main_setmainfont:nn
      { \c_@@_name_gyrefeature_clist } { texgyretermes }
    \__fontspec_main_setsansfont:nn
      { \c_@@_name_gyrefeature_clist } { texgyreheros }
    \__fontspec_main_setmonofont:nn
      {
        \c_@@_name_gyrefeature_clist,
        Scale = MatchLowercase,
        Ligatures = CommonOff,
      }
      { texgyrecursor }
  }
%    \end{macrocode}
% \end{macro}
%
% \begin{macro}{\@@_loadfont_latin_fandol:}
% 为兼容 \pkg{ctex} 做出的名称改变。
%    \begin{macrocode}
\cs_new_eq:NN \@@_loadfont_latin_fandol: \@@_loadfont_latin_gyre:
%    \end{macrocode}
% \end{macro}
%
% \begin{variable}{\c_@@_name_gyrefeature_clist}
% 用于 \pkg{fontspec} 的 \TeX Gyre 字体特性列表。
%    \begin{macrocode}
\clist_const:Nn \c_@@_name_gyrefeature_clist
  {
    Extension = .otf,
    UprightFont = *-regular,
    BoldFont = *-bold,
    ItalicFont = *-italic,
    BoldItalicFont = *-bolditalic,
  }
%    \end{macrocode}
% \end{variable}
%
%
% \subsubsection{定义中文字库}
%
% \begin{macro}{\@@_hide_no_script_msg:}
% 隐藏 |does not contain script "CJK"| 警告。
%    \begin{macrocode}
\cs_new:Npn \@@_hide_no_script_msg:
  { \msg_redirect_name:nnn { fontspec } { no-script } { info } }
%    \end{macrocode}
% \end{macro}
%
% \begin{macro}{\@@_loadfont_cjk_win:N, \@@_loadfont_cjk_win:}
% Windows 中文字体。
%    \begin{macrocode}
\cs_new:Npn \@@_loadfont_cjk_win:N #1
  {
    \setCJKmainfont { SimSun } [ #1, ItalicFont = KaiTi ]
    \setCJKsansfont { SimHei } [#1]
    \setCJKmonofont { FangSong } [#1]
    \setCJKfamilyfont { zhsong } { SimSun } [#1]
    \setCJKfamilyfont { zhhei } { SimHei } [#1]
    \setCJKfamilyfont { zhfs } { FangSong } [#1]
    \setCJKfamilyfont { zhkai } { KaiTi } [#1]
    \@@_define_stzhongs:nn { STZhongsong } {#1}
  }
\cs_new:Npn \@@_loadfont_cjk_win:
  { \@@_loadfont_cjk_win:N \c_@@_name_fakebold_tl }
%    \end{macrocode}
% \end{macro}
%
% \begin{macro}{\@@_loadfont_cjk_mac:}
% macOS 字体。
% ^^A TODO: 修复 macOS 字体支持，实现开箱即用。
%    \begin{macrocode}
\cs_new_protected:Npn \@@_loadfont_cjk_mac:
  {
    \@@_hide_no_script_msg:
    \setCJKmainfont { Songti~SC~Light }
      [
        BoldFont = Songti~SC~Bold,
        ItalicFont = Kaiti~SC,
        BoldItalicFont = Kaiti~SC~Bold,
      ]
    \setCJKsansfont { Heiti~SC~Light } [ BoldFont = Heiti~SC~Medium ]
    \setCJKmonofont { STFangsong }
    \setCJKfamilyfont { zhsong } { Songti~SC~Light } [ BoldFont = Songti~SC~Bold ]
    \setCJKfamilyfont { zhhei } { Heiti~SC~Light } [ BoldFont = Heiti~SC~Medium ]
    \setCJKfamilyfont { zhfs } { STFangsong }
    \setCJKfamilyfont { zhkai } { Kaiti~SC } [ BoldFont = Kaiti~SC~Bold ]
  }
%    \end{macrocode}
% \end{macro}
%
% \begin{macro}{\@@_loadfont_cjk_macoffice:n, \@@_loadfont_cjk_macoffice:}
% mac Office 中文字体。
%    \begin{macrocode}
\cs_new:Npn \@@_loadfont_cjk_macoffice:n #1
  {
    \setCJKmainfont { Simsun.ttc } [ ItalicFont = Kaiti.ttf, #1 ]
    \setCJKsansfont { SimHei.ttf } [#1]
    \setCJKmonofont { Fangsong.ttf } [#1]
    \setCJKfamilyfont { zhsong } { Simsun.ttc } [#1]
    \setCJKfamilyfont { zhhei } { SimHei.ttf } [#1]
    \setCJKfamilyfont { zhfs } { Fangsong.ttf } [#1]
    \setCJKfamilyfont { zhkai } { Kaiti.ttf } [#1]
    \@@_define_stzhongs:nn { \c__hust_path_macoffice_tl STZHONGS.ttf } {#1}
  }
\cs_new:Npn \@@_loadfont_cjk_macoffice:
  {
    \@@_loadfont_cjk_macoffice:n
      { Path = \c_@@_path_macoffice_tl, \c_@@_name_fakebold_tl }
  }
%    \end{macrocode}
% \end{macro}
%
% \begin{macro}{\@@_loadfont_cjk_fandol:}
% Fandol 字体。
%    \begin{macrocode}
\cs_new_protected:Npn \@@_loadfont_cjk_fandol:
  {
    \@@_hide_no_script_msg:
    \setCJKmainfont { FandolSong-Regular }
      [
        Extension = .otf,
        BoldFont = FandolSong-Bold,
        ItalicFont = FandolKai-Regular,
      ]
    \setCJKsansfont { FandolHei-Regular }
      [
        Extension = .otf,
        BoldFont = FandolHei-Bold,
      ]
    \setCJKmonofont { FandolFang-Regular }
      [ Extension = .otf ]
    \setCJKfamilyfont { zhsong } { FandolSong-Regular }
      [
        Extension = .otf,
        BoldFont  = FandolSong-Bold
      ]
    \setCJKfamilyfont { zhhei } { FandolHei-Regular }
      [
        Extension = .otf,
        BoldFont = FandolHei-Bold,
      ]
    \setCJKfamilyfont { zhfs } { FandolFang-Regular }
      [ Extension = .otf ]
    \setCJKfamilyfont { zhkai } { FandolKai-Regular }
      [ Extension = .otf, \c_@@_name_fakebold_tl ]
  }
%    \end{macrocode}
% \end{macro}
%
% \begin{macro}{\@@_loadfont_cjk_founder:}
% 方正字库（简繁扩展）
%    \begin{macrocode}
\cs_new_protected:Npn \@@_loadfont_cjk_founder:
  {
%    \end{macrocode}
% 调整方正字体括号位置。
% \footnote{\XeTeX 的调整方法来自 \url{https://www.zhihu.com/question/46241367/answer/101660183}。}
%    \begin{macrocode}
    \sys_if_engine_xetex:T
      { \xeCJKEditPunctStyle { quanjiao } { optimize-kerning = true } }
    \sys_if_engine_luatex:T
      { \defaultCJKfontfeatures { JFM = { zh_CN / { quanjiao, fzpr } } } }
    \setCJKmainfont { FZShuSong-Z01 }
      [ BoldFont = FZXiaoBiaoSong-B05, ItalicFont = FZKai-Z03 ]
    \setCJKsansfont { FZXiHeiI-Z08 } [ BoldFont = FZHei-B01 ]
    \setCJKmonofont { FZFangSong-Z02 }
    \setCJKfamilyfont { zhsong } { FZShuSong-Z01 }
      [ BoldFont = FZXiaoBiaoSong-B05 ]
    \setCJKfamilyfont { zhhei } { FZHei-B01 }
      [ \c_@@_name_fakebold_tl ]
    \setCJKfamilyfont { zhkai } { FZKai-Z03 }
      [ \c_@@_name_fakebold_tl ]
    \setCJKfamilyfont { zhfs } { FZFangSong-Z02 }
    \defaultCJKfontfeatures { }
  }
%    \end{macrocode}
% \end{macro}
%
% \begin{macro}{\@@_loadfont_cjk_noto:}
% Noto 思源字体。
%    \begin{macrocode}
\cs_new_protected:Npn \@@_loadfont_cjk_noto:
  {
    \setCJKmainfont [ \c_@@_name_notofeature_clist ]
      { NotoSerifCJKsc }
    \setCJKsansfont [ \c_@@_name_notofeature_clist ]
      { NotoSansCJKsc  }
    \setCJKmonofont { Noto~Sans~Mono~CJK~SC }
    \setCJKfamilyfont { zhsong } { Noto~Serif~CJK~SC }
    \setCJKfamilyfont { zhhei } { Noto~Sans~CJK~SC }
    \setCJKfamilyfont { zhfs } { FZFangSong-Z02 }
    \setCJKfamilyfont { zhkai } { FZKai-Z03 }
      [ \c_@@_name_fakebold_tl ]
  }
%    \end{macrocode}
% \end{macro}
%
% \begin{macro}{\@@_loadfont_cjk_source:}
%    \begin{macrocode}
\cs_new_protected:Npn \@@_loadfont_cjk_source:
  {
    \setCJKmainfont [ \c_@@_name_notofeature_clist ]
      { SourceHanSerifSC }
    \setCJKsansfont [ \c_@@_name_notofeature_clist ]
      { SourceHanSansSC  }
    \setCJKmonofont { FZFangSong-Z02 }
    \setCJKfamilyfont { zhsong } { Source~Han~Serif~SC }
    \setCJKfamilyfont { zhhei } { Source~Han~Sans~SC }
    \setCJKfamilyfont { zhfs } { FZFangSong-Z02 }
    \setCJKfamilyfont { zhkai } { FZKai-Z03 }
      [ \c_@@_name_fakebold_tl ]
  }
%    \end{macrocode}
% \end{macro}
%
% \begin{variable}{\c_@@_name_fakebold_tl}
%    \begin{macrocode}
\tl_const:Nn \c_@@_name_fakebold_tl { AutoFakeBold = 2.17 }
%    \end{macrocode}
% \end{variable}
%
% \begin{variable}{\c_@@_name_notofeature_clist}
% 用于 \pkg{fontspec} 的思源字体特性列表。
%    \begin{macrocode}
\clist_const:Nn \c_@@_name_notofeature_clist
  {
    Extension = .otf,
    UprightFont = *-Regular,
    BoldFont = *-Bold,
    ItalicFont = *-Regular,
    BoldItalicFont = *-Bold,
    ItalicFeatures = FakeSlant,
    BoldItalicFeatures = FakeSlant,
  }
%    \end{macrocode}
% \end{variable}
%
% \begin{variable}{\l_@@_name_stzhongsopts_tl}
% 华文中宋的字体特性列表，默认为仅使用伪粗体。
%    \begin{macrocode}
\tl_new:N \l_@@_name_stzhongsopts_tl
\tl_set_eq:NN \l_@@_name_stzhongsopts_tl \c_@@_name_fakebold_tl
%    \end{macrocode}
% \end{variable}
%
% \begin{macro}{\@@_define_stzhongs:nn}
% 中文封面额外需要的华文中宋。
%    \begin{macrocode}
\tl_new:N \l_@@_name_stzhongsfile_tl
\cs_new:Npn \@@_define_stzhongs:nn #1#2
  {
    \tl_set:Nn \l_@@_name_stzhongsfile_tl {#1}
    \tl_set:Nn \l_@@_name_stzhongsopts_tl {#2}
  }
%    \end{macrocode}
% \end{macro}
%
% \begin{macro}{\@@_loadfont_stzhongs:n}
% \begin{macro}{\@@_loadfont_stzhongs:V, \@@_loadfont_stzhongs:, \@@_stzhongs:}
% 研究生封面额外需要的华文中宋。
%    \begin{macrocode}
\cs_new:Npn \@@_loadfont_stzhongs:
  {
    \@@_loadfont_stzhongs:V \l_@@_name_stzhongsfile_tl
  }
\cs_new_protected:Npn \@@_loadfont_stzhongs:n #1
  {
    \fontspec_font_if_exist:nTF { #1 }
      {
        \newCJKfontfamily \@@_stzhongs: { #1 }
          [ \l_@@_name_stzhongsopts_tl ]
      }
      {
        \bool_if:NTF \g_@@_font_dir_bool
          {
            \newCJKfontfamily \@@_stzhongs: { #1 }
              [ Path = \g_@@_font_dir_tl, \c_@@_name_fakebold_tl ]
          }
          {
            \cs_set_eq:NN \@@_stzhongs: \rmfamily
            \@@_warning:n { missing-stzhongs }
          }
      }
  }
\@@_msg_new:nn { missing-stzhongs }
  {
    "STZHONGS.TTF"~ is~ not~ found~ in~ your~ system.\\
    The~ font~ is~ normally~ distributed~ with~ MS~ Windows.~
    I~ have~ used~ songti~ for~ substitution.
  }
\cs_generate_variant:Nn \@@_loadfont_stzhongs:n { V }
%    \end{macrocode}
% \end{macro}
% \end{macro}
%
%
% \subsubsection{定义数学字库}
%
% \begin{macro}{\@@_define_math_font:nn}
% 批量定义数学字体配置。
%    \begin{macrocode}
\cs_new:Npn \@@_define_math_font:nn #1#2
  {
    \cs_new:cpn { @@_loadfont_math_ #1 : }
      { \__um_setmathfont:nn { } {#2} }
  }
%    \end{macrocode}
% \end{macro}
%
% \begin{macro}
% {
%   \@@_loadfont_math_asana:,
%   \@@_loadfont_math_fira:,
%   \@@_loadfont_math_garamond:,
%   \@@_loadfont_math_lm:,
%   \@@_loadfont_math_libertinus:,
%   \@@_loadfont_math_stix:,
%   \@@_loadfont_math_bonum:,
%   \@@_loadfont_math_dejavu:,
%   \@@_loadfont_math_pagella:,
%   \@@_loadfont_math_schola:,
%   \@@_loadfont_math_termes:,
% }
% 批量定义若干数学字体的载入命令。
%    \begin{macrocode}
\clist_map_inline:nn
  {
    { asana } { Asana-Math.otf },
    { fira } { FiraMath-Regular.otf },
    { garamond } { Garamond-Math.otf },
    { lm } { latinmodern-math.otf },
    { libertinus } { LibertinusMath-Regular.otf },
    { stix } { STIXMath-Regular.otf },
    { bonum } { texgyrebonum-math.otf },
    { dejavu } { texgyredejavu-math.otf },
    { pagella } { texgyrepagella-math.otf },
    { schola } { texgyreschola-math.otf },
    { termes } { texgyretermes-math.otf },
  }
  { \@@_define_math_font:nn #1 }
%    \end{macrocode}
% \end{macro}
%
% \begin{macro}{\@@_loadfont_math_cambria:}
% Cambria Math 字体配置。
%    \begin{macrocode}
\cs_new:Npn \@@_loadfont_math_cambria:
  {
    \bool_if:NTF \g_@@_font_dir_bool
      {
        \__um_setmathfont:nn
          { Path = \g_@@_font_dir_tl/, FontIndex = 1 }
          { cambria.ttc }
      }
      { \__um_setmathfont:nn { } { Cambria~Math } }
  }
%    \end{macrocode}
% \end{macro}
%
% \begin{macro}{\@@_loadfont_math_xits:}
% XITS Math 字体。
%    \begin{macrocode}
\cs_new:Npn \@@_loadfont_math_xits:
  {
    \bool_if:NTF \g_@@_opt_math_int_bool
      { \tl_set:Nn \l_@@_tmpa_tl { 8 } }
      { \tl_clear:N \l_@@_tmpa_tl }
    \__um_setmathfont:nn
      {
        Extension = .otf,
        StylisticSet = \l_@@_tmpa_tl,
        BoldFont = XITSMath-Bold,
      }
      { XITSMath-Regular }
    \__um_setmathfont:nn
      {
        Extension = .otf,
        StylisticSet = 1,
        range = {cal, bfcal},
      }
      { XITSMath-Regular }
  }
%    \end{macrocode}
% \end{macro}
%
% \begin{macro}{\@@_loadfont_math_newcm:}
% New Computer Modern Math 字体。
%    \begin{macrocode}
\cs_new:Npn \@@_loadfont_math_newcm:
  {
    \bool_if:NTF \g_@@_opt_math_int_bool
      { \tl_set:Nn \l_@@_tmpa_tl { 2 } }
      { \tl_clear:N \l_@@_tmpa_tl }
    \__um_setmathfont:nn
      {
        Extension = .otf,
        StylisticSet = \l_@@_tmpa_tl,
      }
      { NewCMMath-Book }
    \__um_setmathfont:nn
      {
        Extension = .otf,
        StylisticSet = 1,
        range = {scr,bfscr},
      }
      { NewCMMath-Book }
    \__fontspec_main_setmathrm:nn
      {
        Extension = .otf,
        UprightFont = *-Book,
        BoldFont = *-Bold,
        ItalicFont = *-BookItalic,
        BoldItalicFont = *-BoldItalic,
      }
      { NewCM10 }
    \__fontspec_main_setmathsf:nn
      {
        Extension = .otf,
        UprightFont = *-Book,
        BoldFont = *-Bold,
        ItalicFont = *-BookOblique,
        BoldItalicFont = *-BoldOblique,
      }
      { NewCMSans10 }
    \__fontspec_main_setmathtt:nn
      {
        Extension = .otf,
        UprightFont = *-Book,
        ItalicFont = *-BookItalic,
        BoldFont = *-Bold,
        BoldItalicFont = *-BoldOblique,
      }
      { NewCMMono10 }
  }
%    \end{macrocode}
% \end{macro}
%
% \begin{macro}{\@@_loadfont_math_none:}
% 不进行数学字体配置。
%    \begin{macrocode}
\cs_set_eq:NN \@@_loadfont_math_none: \tex_relax:D
%    \end{macrocode}
% \end{macro}
%
%
% \subsubsection{载入指定字库}
%
% \begin{macro}{\@@_loadfont:}
% 载入字体命令。
%    \begin{macrocode}
\cs_new_protected:Npn \@@_loadfont:
  {
    \use:c { @@_loadfont_latin_ \g_@@_font_latin_tl : }
    \use:c { @@_loadfont_cjk_ \g_@@_font_cjk_tl : }
%    \end{macrocode}
% 自行定义 \pkg{ctex} 中的四类字体命令。
%    \begin{macrocode}
    \NewDocumentCommand \songti { } { \CJKfamily { zhsong } }
    \NewDocumentCommand \heiti { } { \CJKfamily { zhhei } }
    \NewDocumentCommand \fangsong { } { \CJKfamily { zhfs } }
    \NewDocumentCommand \kaishu { } { \CJKfamily { zhkai } }
%    \end{macrocode}
% \begin{macro}{\bigger}
% 重定义字号命令。
%    \begin{macrocode}
    \NewDocumentCommand \bigger { } { \ctex_zihao:n { 4 } }
  }
%    \end{macrocode}
% \end{macro}
% \end{macro}
%
% 载入设置的字体。^^A 为了吸收导言区的设置，放在其后载入。
%    \begin{macrocode}
% \BeforeBeginEnvironment { document } { \@@_loadfont: }
\@@_loadfont:
%    \end{macrocode}
%
% 将华文中宋放在文档类最后载入。
%    \begin{macrocode}
\ctex_at_end_preamble:n { \@@_loadfont_stzhongs: }
%    \end{macrocode}
%
% \subsection{页眉页脚}
%
% 清除默认页眉页脚格式。
%    \begin{macrocode}
\fancyhf {}
%    \end{macrocode}
%
% 设置页眉页脚颜色。
% \changes{v2.0}{2025-02-21}{使用 \pkg{l3color}，不再依赖 \pkg{xcolor}}
%    \begin{macrocode}
\color_set:nn { headfoot } { red }
% \providecolor { headfoot } { rgb } { 1, 0, 0 }
%    \end{macrocode}
%
% \begin{macro}{\@@_draw_hf_hline:n}
% 绘制页眉页脚横线。
% \begin{arguments}
%   \item 页眉页脚线的高度（粗细）
% \end{arguments}
%    \begin{macrocode}
\cs_new_protected:Npn \@@_draw_hf_hline:n #1
  {
    { \color_select:n { headfoot } \tex_hrule:D height #1 }
  }
%    \end{macrocode}
% \end{macro}
%
% 设置页眉格式。仅在非盲审模式~\key{anonymous} 下显示学校名称。
%    \begin{macrocode}
\bool_if:NF \g_@@_anonymous_bool
  {
    \chead
      {
        \zihao{3} \ziju{0.7} \bfseries \kaishu
        \color_group_begin:
        \color_select:n { headfoot }
        \c_@@_header_tl
        \color_group_end:
      }
  }
\tl_set:Nn { \headrulewidth } { 1 pt }
\tl_set:Nn { \headrule }
  {
    \skip_vertical:n { 6 pt }
    \@@_draw_hf_hline:n \headrulewidth
    \skip_vertical:n \headrulewidth
    \@@_draw_hf_hline:n \headrulewidth
  }
%    \end{macrocode}
%
% 设置页脚格式。
%    \begin{macrocode}
\cfoot { \zihao{-5} \thepage }
\tl_set:Nn { \footrulewidth } { 0.5 pt }
\tl_set:Nn { \footrule }
  {
    \@@_draw_hf_hline:n \footrulewidth
  }
%    \end{macrocode}
%
% \pkg{ctex} 宏包使用 |heading| 选项后，会把页面格式设置为 |headings|。
% 因此必须在 \pkg{ctex} 调用之后重新设置 \tn{pagestyle} 为 |fancy|。
%    \begin{macrocode}
\pagestyle { fancy }
%    \end{macrocode}
%
%
% \section{封面样式}
%
% \subsection{信息录入}
%
% \begin{variable}{
%   \g_@@_info_title_tl,
%   \g_@@_info_degree_tl,
%   \g_@@_info_author_tl,
%   \g_@@_info_major_tl,
%   \g_@@_info_supervisor_tl,
%   \g_@@_info_date_tl,
%   \g_@@_info_clc_tl,
%   \g_@@_info_student_id_tl,
%   \g_@@_info_school_id_tl,
%   \g_@@_info_secret_tl,
% }
% 声明中文封面信息。
%    \begin{macrocode}
\tl_new:N \g_@@_info_title_tl
\tl_new:N \g_@@_info_degree_type_tl
\tl_new:N \g_@@_info_author_tl
\tl_new:N \g_@@_info_major_tl
\tl_new:N \g_@@_info_supervisor_tl
\tl_new:N \g_@@_info_date_tl
\tl_new:N \g_@@_info_clc_tl
\tl_new:N \g_@@_info_student_id_tl
\tl_new:N \g_@@_info_school_id_tl
\tl_new:N \g_@@_info_secret_tl
%    \end{macrocode}
% \end{variable}
%
% \begin{variable}{
%   \g_@@_info_title_en_tl,
%   \g_@@_info_author_en_tl,
%   \g_@@_info_major_en_tl,
%   \g_@@_info_supervisor_en_tl,
%   \g_@@_info_date_en_tl,
%   \g_@@_info_degree_en_tl,
% }
% 声明英文封面信息。
%    \begin{macrocode}
\tl_new:N \g_@@_info_title_en_tl
\tl_new:N \g_@@_info_author_en_tl
\tl_new:N \g_@@_info_major_en_tl
\tl_new:N \g_@@_info_supervisor_en_tl
\tl_new:N \g_@@_info_date_en_tl
\tl_new:N \g_@@_info_degree_en_tl
%    \end{macrocode}
% \end{variable}
%
% \begin{variable}{\l_@@_info_title_sanitized_tl,\l_@@_info_title_en_sanitized_tl}
% 处理后的论文标题。
%    \begin{macrocode}
\tl_new:N \l_@@_info_title_sanitized_tl
\tl_new:N \l_@@_info_title_en_sanitized_tl
%    \end{macrocode}
% \end{variable}
%
% \begin{macro}[do-not-index={\1, \2}]{\@@_set_title:nn}
% 设置论文标题，并删除其中的 |\\| 和汉字间的空格。
%    \begin{macrocode}
\cs_new_protected:Npn \@@_set_title:nn #1#2
  {
    \tl_gset:cn { g_@@_info_ #1 _tl } { #2 }
    \tl_set:cn { l_@@_info_ #1 _sanitized_tl } { #2 }
    \tl_replace_all:cnn { l_@@_info_ #1 _sanitized_tl } { \\ } { ~ }
    \regex_replace_all:nnc
      { \s+ } { \ }
      { l_@@_info_ #1 _sanitized_tl }
    \regex_replace_all:nnc
      { ([\x{4e00}-\x{9fff}]) \s+ ([\x{4e00}-\x{9fff}]) } { \1 \2 }
      { l_@@_info_ #1 _sanitized_tl }
  }
%    \end{macrocode}
% \end{macro}
%
% \begin{macro}{\@@_set_date:n}
% 设置论文提交日期。
%    \begin{macrocode}
\cs_new:Npn \@@_set_date:n #1
  {
    \@@_set_date_aux:w #1 \q_stop
  }
\cs_new_protected:Npn \@@_set_date_aux:w #1 - #2 - #3 \q_stop
  {
    \tl_gset:Nn \g_@@_info_date_tl { {#1} 年 {#2} 月 {#3} 日 }
    \tl_gset:Nn \g_@@_info_date_en_tl
      { \clist_item:Nn \c_@@_month_clist {#2},~ {#1} }
  }
\clist_const:Nn \c_@@_month_clist
  {
    January, February, March, April, May, June,
    July, August, September, October, November, December,
  }
%    \end{macrocode}
% \end{macro}
%
% 定义 |hust/info| 键值类。
%    \begin{macrocode}
\keys_define:nn { hust / info }
  {
%    \end{macrocode}
%
% \begin{macro}[module=info/]{info/title, info/title*}
% 论文题目。以下带星号的项目均表示相应的英文字段。
%    \begin{macrocode}
    title .tl_gset:N = \g_@@_info_title_tl,
    title* .tl_gset:N = \g_@@_info_title_en_tl,
%    title .code:n = { \@@_set_title:nn { title } { #1 } },
%    title* .code:n = { \@@_set_title:nn { title_en } { #1 } },
%    \end{macrocode}
% \end{macro}
%
% \begin{macro}[module=info/]{info/degree, info/degree*}
% 学位类型（中文）或学科（英文）。
% 中文类型~\key[info]{degree} 分为学术型~\opt{academic}，和专业型~\opt{professional}；
% 英文学科~\key[info]{degree*}\footnote{为保持命名一致性，学科 degree discipline 在此处仍使用键名 |degree*|} 包括文学~\opt{Arts}，理学~\opt{Science}，教育学~\opt{Education}，工学~\opt{Engineering}，法学~\opt{Laws}，医学~\opt{Medicine}，以及其他~\opt{XXX}。
%    \begin{macrocode}
    degree .choices:nn =
      { academic, professional }
      {
        \tl_gset_eq:NN \g_@@_info_degree_tl \l_keys_choice_tl
      },
    degree .default:n = { academic },
    degree* .tl_gset:N = \g_@@_info_degree_en_tl,
%    \end{macrocode}
% \end{macro}
%
% \begin{macro}[module=info/]{info/author, info/author*}
% 学位申请人姓名。
%    \begin{macrocode}
    author .tl_gset:N = \g_@@_info_author_tl,
    author* .tl_gset:N = \g_@@_info_author_en_tl,
%    \end{macrocode}
% \end{macro}
%
% \begin{macro}[module=info/]{info/major, info/major*}
% 专业名称。
%    \begin{macrocode}
    major .tl_gset:N = \g_@@_info_major_tl,
    major* .tl_gset:N = \g_@@_info_major_en_tl,
%    \end{macrocode}
% \end{macro}
%
% \begin{macro}[module=info/]{info/supervisor, info/supervisor*}
% 指导老师列表。
%    \begin{macrocode}
    supervisor .tl_gset:N = \g_@@_info_supervisor_tl,
    supervisor* .tl_gset:N = \g_@@_info_supervisor_en_tl,
%    \end{macrocode}
% \end{macro}
%
% \begin{macro}[module=info/]{info/date}
% 论文提交日期。
%    \begin{macrocode}
    date .code:n = { \@@_set_date:n { #1 } },
%    \end{macrocode}
% \end{macro}
%
% \begin{macro}[module=info/]{info/clc}
% 中图分类号。
%    \begin{macrocode}
    clc .tl_gset:N = \g_@@_info_clc_tl,
%    \end{macrocode}
% \end{macro}
%
% \begin{macro}[module=info/]{info/student-id, info/school-id}
% 学号和学校代码。
%    \begin{macrocode}
    student-id .tl_gset:N = \g_@@_info_student_id_tl,
    school-id .tl_gset:N = \g_@@_info_school_id_tl,
%    \end{macrocode}
% \end{macro}
%
% \begin{macro}[module=info/]{info/secret}
% 密级。
%    \begin{macrocode}
    secret .tl_gset:N = \g_@@_info_secret_tl,
%    \end{macrocode}
% \end{macro}
%
%    \begin{macrocode}
  }
%    \end{macrocode}
%
% \subsubsection{定义内部函数}
%
% \begin{macro}{\@@_spread_box:nn, \@@_spread_box:no}
% 分散对齐的水平盒子。
% \begin{arguments}
%   \item 宽度
%   \item 内容
% \end{arguments}
% 利用 \cs{tl_map_inline:nn} 在字符间插入 \tn{hfil}；紧随其后的 \tn{unskip}
% 将会去掉最后一个 \tn{hfil}。见 \url{https://tex.stackexchange.com/q/169689}。
%    \begin{macrocode}
\cs_new_protected:Npn \@@_spread_box:nn #1#2
  {
    \mode_leave_vertical:
    \hbox_to_wd:nn {#1} { \tl_map_inline:nn {#2} { ##1 \hfil } \unskip }
  }
\cs_generate_variant:Nn \@@_spread_box:nn { no }
%    \end{macrocode}
% \end{macro}
%
% \begin{macro}{\@@_center_box:nn, \@@_center_box:Vn}
% 居中对齐的水平盒子。
%    \begin{macrocode}
\cs_new_protected:Npn \@@_center_box:nn #1#2
  {
    \mode_leave_vertical:
    \hbox_to_wd:nn {#1} { \tex_hfil:D #2 \tex_hfil:D }
  }
\cs_generate_variant:Nn \@@_center_box:nn { Vn }
%    \end{macrocode}
% \end{macro}
%
% \begin{macro}{\@@_fixed_width_box:nn}
% 限宽盒子（允许换行）。
%    \begin{macrocode}
\cs_new:Npn \@@_fixed_width_box:nn #1#2
  { \parbox {#1} {#2} }
%    \end{macrocode}
% \end{macro}
%
% \begin{macro}{\@@_fixed_width_center_box:nn}
% 居中对齐的限宽盒子（允许换行）。
%    \begin{macrocode}
\cs_new:Npn \@@_fixed_width_center_box:nn #1#2
  { \parbox {#1} { \centering #2 } }
%    \end{macrocode}
% \end{macro}
%
% \begin{macro}{\@@_get_text_width:Nn, \@@_get_text_width:NV}
% 获取文本宽度，并存入 |dim| 型变量。
% \begin{arguments}
%   \item |dim| 型变量
%   \item 内容
% \end{arguments}
%    \begin{macrocode}
\cs_new:Npn \@@_get_text_width:Nn #1#2
  {
    \hbox_set:Nn \l_@@_tmpa_box {#2}
    \dim_set:Nn #1 { \box_wd:N \l_@@_tmpa_box }
  }
\cs_generate_variant:Nn \@@_get_text_width:Nn { NV }
%    \end{macrocode}
% \end{macro}
%
% \begin{macro}[deprecated]{\@@_blank_underline:n}
% 下划线占位符。|#1|: 长度。
%    \begin{macrocode}
\cs_new:Npn \@@_blank_underline:n #1
  { \rule [ -0.5 ex ] {#1} { 0.4 pt } }
%    \end{macrocode}
% \end{macro}
%
% \subsection{封面模板}
%
% \begin{macro}[int, module=hust/component/]{component/}
% 封面组件。
% \begin{arguments}
%   \item[\opt{content}] 组件主体内容
%   \item[\opt{format}] 组件字体/段落格式
%   \item[\opt{height}] 组件固定高度，不设置则自动计算。主要用于行数不确定，需要固定高度的情况
%   \item[\opt{top-skip}] 顶部间距（建议优先配置底部间距，顶部间距仅在特殊情况使用）
%   \item[\opt{bottom-skip}] 底部间距
%   \item[\opt{align}] 对齐方式，可选值为靠左对齐~\opt{left}，靠右对齐~\opt{right}，居中对齐~\opt{center}，以及两端对齐~\opt{justify}。默认为居中对齐
% \end{arguments}
%    \begin{macrocode}
\NewTemplateType { hust / component } { 0 }
\DeclareTemplateInterface { hust / component } { plain } { 0 }
{
  content : tokenlist = \c_empty_tl,
  format : tokenlist = \c_empty_tl,
  width: skip = \c_zero_skip,
  height: skip = \c_zero_skip,
  top-skip : skip = \c_zero_skip,
  bottom-skip : skip = \c_zero_skip,
  align : choice { left, right, center, justify } = center,
}
\DeclareTemplateCode { hust / component } { plain } { 0 }
{
  content = \l_@@_component_content_tl,
  format = \l_@@_component_format_tl,
  width = \l_@@_component_width_skip,
  height = \l_@@_component_height_skip,
  top-skip = \l_@@_component_top_skip,
  bottom-skip = \l_@@_component_bottom_skip,
  align =
    {
      left =
        \cs_set_eq:NN \l_@@_component_align: \raggedright,
      right =
        \cs_set_eq:NN \l_@@_component_align: \raggedleft,
      center =
        \cs_set_eq:NN \l_@@_component_align: \centering,
      justify =
        \cs_set_eq:NN \l_@@_component_align: \prg_do_nothing:,
    }
}
{
  \AssignTemplateKeys
  \dim_compare:nNnTF { \l_@@_component_height_skip } > { \c_zero_skip }
    {
      \box_clear_new:N \l_@@_component_box
      \vbox_set_to_ht:Nnn \l_@@_component_box { \l_@@_component_height_skip }
        {
          \skip_vertical:N \l_@@_component_top_skip
          \l_@@_component_align:
          \l_@@_component_format_tl
          \l_@@_component_content_tl
          \tex_vfill:D
        }
      \box_use:N \l_@@_component_box
    }
    {
      \skip_vertical:N \l_@@_component_top_skip
      \group_begin:
      \l_@@_component_align:
      \l_@@_component_format_tl
      \l_@@_component_content_tl
      \tex_par:D
      \group_end:
    }
  \skip_vertical:N \l_@@_component_bottom_skip
  \skip_vertical:N \c_zero_skip
}
%    \end{macrocode}
% \end{macro}
%
% \begin{macro}[int, module=hust/page]{page/}
% 页面模板。
% \begin{arguments}
%   \item[\opt{prefix}] 组件共用前缀
%   \item[\opt{components}] 组件列表
%   \item[\opt{style}] 页面样式 （可使用 \tn{thispagestyle} 配置）
%   \item[\opt{format}] 页面字体/段落格式
%   \item[\opt{top-skip}] 顶部间距
% \end{arguments}
%    \begin{macrocode}
\NewTemplateType { hust / page } { 0 }
\DeclareTemplateInterface { hust / page } { title } { 0 }
  {
    prefix : tokenlist,
    components : commalist,
    style : tokenlist = empty,
    format : tokenlist = \linespread { } \selectfont,
    top-skip : skip = \c_zero_skip,
  }
\DeclareTemplateCode { hust / page } { title } { 0 }
  {
    prefix = \l_@@_page_prefix_tl,
    components = \l_@@_page_components_clist,
    style = \l_@@_page_style_tl,
    format = \l_@@_page_format_tl,
    top-skip = \l_@@_page_top_skip,
  }
  {
    \AssignTemplateKeys
    \skip_vertical:N \l_@@_page_top_skip
    \group_begin:
      \exp_args:No \thispagestyle { \l_@@_page_style_tl }
      \l_@@_page_format_tl
      \clist_map_inline:Nn \l_@@_page_components_clist
        { \UseInstance { hust / component } { \l_@@_page_prefix_tl / ##1 } }
    \group_end:
    \clearpage
  }
%    \end{macrocode}
% \end{macro}
%
% \begin{macro}{
%   \@@_declare_component_instance:nnn,
%   \@@_declare_page_instance:nn,
% }
% 声明模板实例的辅助函数。
%    \begin{macrocode}
\cs_new_protected:Npn \@@_declare_component_instance:nnn #1#2#3
  { \DeclareInstance { hust / component } { #1/#2 } { plain } {#3} }
\cs_new_protected:Npn \@@_declare_page_instance:nn #1#2
  {
    \DeclareInstance { hust / page } {#1} { title }
    {
      prefix = #1,
      #2,
    }
  }
%    \end{macrocode}
% \end{macro}
%
% \subsection{封面组件}
%
% \subsubsection{中文封面}
%
% \begin{macro}[int, module=hust/component/]{title zh/identifiers}
% 分类号，学号，学校代码，密级。
%    \begin{macrocode}
\cs_new:Npn \@@_title_zh_ids_aux:nn #1#2
  {
    \tl_use:c { c_@@_name_ #2 _tl }
    \@@_ul_to_wd:nnn { 1 pt } { #1 }
      { \tl_use:c { g_@@_info_ #2 _tl } }
  }
\@@_declare_component_instance:nnn { title zh } { identifiers }
  {
    content =
      {
        \@@_title_zh_ids_aux:nn { 6em } { clc } % 3 em
        \skip_horizontal:n { \linewidth - 19em }
        \@@_title_zh_ids_aux:nn { 8em } { student_id } % 4 em
        \tex_par:D
        \@@_title_zh_ids_aux:nn { 5em } { school_id } % 4 em
        \skip_horizontal:n { \linewidth - 19em }
        \@@_title_zh_ids_aux:nn { 8em } { secret } % 4 em
      },
    format =
      {
        \zihao{-4} \bfseries \rmfamily \songti
        \@@_line_spread:n { 1.5 }
      },
    bottom-skip = 2.5 cm,
  }
%    \end{macrocode}
% \end{macro}
%
% \begin{macro}[int, module=hust/component/]{title zh/logo}
% 校名图片。
%    \begin{macrocode}
\@@_declare_component_instance:nnn { title zh } { logo }
  {
    content =
      {
        \includegraphics [ height = 1.4 cm ] { hust-title.pdf }
      },
    bottom-skip = 0.5 cm,
  }
%    \end{macrocode}
% \end{macro}
%
% \begin{macro}[int, module=hust/component/]{title zh/level}
% 学位级别。
%    \begin{macrocode}
\__ctex_save_font_size:nn { a } { 45 bp } % 大初号
\@@_declare_component_instance:nnn { title zh } { level }
  {
    content = { \c_@@_name_level_zh_tl },
    format = { \zihao { a } \bfseries \songti \ziju { 0.1 } },
    bottom-skip = 0.5 cm,
  }
%    \end{macrocode}
% \end{macro}
%
% \begin{macro}[int, module=hust/component/]{title zh/degree}
% 学位类型。
%    \begin{macrocode}
\@@_declare_component_instance:nnn { title zh } { degree }
  {
    content =
      {
        （
        \c_@@_name_academic_zh_tl
        \tl_if_eq:NnTF \g_@@_info_degree_tl { academic }
          { \c_@@_checkbox_tl }
          { \c_@@_emptybox_tl }
        \skip_horizontal:n { 2.5 em }
        \c_@@_name_professional_zh_tl
        \tl_if_eq:NnTF \g_@@_info_degree_tl { professional }
          { \c_@@_checkbox_tl }
          { \c_@@_emptybox_tl }
        ）
      },
    format = { \zihao { -3 } \c_@@_name_degree_font_tl },
    bottom-skip = 1.5 cm,
  }
%    \end{macrocode}
% \end{macro}
%
% \begin{macro}[int, module=hust/component/]{title zh/title}
% 论文题目。
%    \begin{macrocode}
\@@_declare_component_instance:nnn { title zh } { title }
  {
    content = { \g_@@_info_title_tl },
    format =
      {
        \zihao { 1 } \bfseries \songti
        \@@_line_spread:n { 1.5 }
      },
    height = 4.8 cm,
    bottom-skip = 1.2 cm,
  }
%    \end{macrocode}
% \end{macro}
%
% \begin{macro}[int, module=hust/component/]{title zh/info}
% 信息表格。
%    \begin{macrocode}
\@@_declare_component_instance:nnn { title zh } { info }
  {
    content =
      {
        \tl_set:Nn \arraystretch { 2 }
        \begin{tabular}{p{5em}@{：\hspace{1em}}l}
          \c_@@_name_author_zh_tl & \g_@@_info_author_tl \\
          \c_@@_name_major_zh_tl & \g_@@_info_major_tl \\
          \c_@@_name_supervisor_zh_tl & \g_@@_info_supervisor_tl \\
          \c_@@_name_date_tl & \g_@@_info_date_tl \\
        \end{tabular}
      },
    format = { \zihao { -3 } \bfseries \rmfamily \songti },
    bottom-skip = 30 pt,
  }
%    \end{macrocode}
% \end{macro}
%
% \subsubsection{英文封面}
%
% \begin{macro}[int, module=hust/component/]{title en/degree}
% 学位级别和学科。
%    \begin{macrocode}
\@@_declare_component_instance:nnn { title en } { degree }
  {
    content =
      {
        \c_@@_name_level_en_tl
        \tl_if_eq:NnTF \g_@@_info_degree_tl { academic }
          { \c_@@_name_academic_en_tl }
          { \c_@@_name_professional_en_tl }
        \g_@@_info_degree_en_tl
      },
    format =
      {
        \zihao { -3 } \bfseries \rmfamily
        \@@_line_spread:n { 1.5 }
      },
    bottom-skip = 5 cm,
  }
%    \end{macrocode}
% \end{macro}
%
% \begin{macro}[int, module=hust/component/]{title en/title}
% 论文题目。
%    \begin{macrocode}
\@@_declare_component_instance:nnn { title en } { title }
  {
    content = { \g_@@_info_title_en_tl },
    format =
      {
        \zihao { -2 } \bfseries \rmfamily
        \@@_line_spread:n { 1.5 }
      },
    height = 3 cm,
    bottom-skip = 1.2 cm,
  }
%    \end{macrocode}
% \end{macro}
%
% \begin{macro}[int, module=hust/component/]{title en/info}
% 信息表格。
%    \begin{macrocode}
\@@_declare_component_instance:nnn { title en } { info }
  {
    content =
      {
        \begin{tabular}{l@{\hspace{0.5em}:\hspace{0.5em}}l}
          \c_@@_name_author_en_tl & \g_@@_info_author_en_tl \\
          \c_@@_name_major_en_tl & \g_@@_info_major_en_tl \\
          \c_@@_name_supervisor_en_tl & \g_@@_info_supervisor_en_tl \\
        \end{tabular}
      },
    format =
      {
        \zihao { -3 } \bfseries \rmfamily
        \@@_line_spread:n { 1.5 }
      },
    bottom-skip = 4 cm,
  }
%    \end{macrocode}
% \end{macro}
%
% \begin{macro}[int, module=hust/component/]{title en/address}
% 学校名称和地址。
%    \begin{macrocode}
\@@_declare_component_instance:nnn { title en } { address }
  {
    content = { \c_@@_address_tl },
    format =
      {
        \zihao { -3 } \bfseries \rmfamily
        \@@_line_spread:n { 1.5 }
      },
  }
%    \end{macrocode}
% \end{macro}
%
% \begin{macro}[int, module=hust/component/]{title en/date}
%    \begin{macrocode}
\@@_declare_component_instance:nnn { title en } { date }
  {
    content = { \tl_use:N \g_@@_info_date_en_tl },
    format =
      {
        \zihao { -3 } \bfseries \rmfamily
        \@@_line_spread:n { 1.5 }
      },
  }
%    \end{macrocode}
% \end{macro}
%
%
% \subsubsection{声明页}
%
% \begin{macro}[int, module=hust/component/]{decl/orig title}
% 独创性声明标题。
%    \begin{macrocode}
\@@_declare_component_instance:nnn { decl } { orig title }
  {
    content = { \c_@@_orig_decl_title_tl },
    format =
      {
        \zihao { 3 } \bfseries \heiti
        \@@_line_spread:n { 1.5 }
      },
    bottom-skip = 1 cm,
  }
%    \end{macrocode}
% \end{macro}
%
% \begin{macro}[int, module=hust/component/]{decl/orig text}
% 独创性声明正文。
%    \begin{macrocode}
\@@_declare_component_instance:nnn { decl } { orig text }
  {
    content = { \c_@@_orig_decl_text_tl },
    format =
      {
        \zihao { -4 } \rmfamily
        \@@_line_spread:n { 1.5 }
      },
    align = justify,
    bottom-skip = 1 cm,
  }
%    \end{macrocode}
% \end{macro}
%
% \begin{macro}[int, module=hust/component/]{decl/orig sign}
% 独创性声明签名。
%    \begin{macrocode}
\@@_declare_component_instance:nnn { decl } { orig sign }
  {
    content =
      {
        \clist_map_inline:Nn \c_@@_orig_decl_sign_clist
          {
            #1 \\
          }
      },
    format =
      {
        \zihao { -4 } \rmfamily
        \@@_line_spread:n { 1.5 }
      },
    align = right,
    bottom-skip = 3.5 cm,
  }
%    \end{macrocode}
% \end{macro}
%
% \begin{macro}[int, module=hust/component/]{decl/auth title}
% 学位论文版权使用授权书标题。
%    \begin{macrocode}
\@@_declare_component_instance:nnn { decl } { auth title }
  {
    content = { \c_@@_auth_decl_title_tl },
    format =
      {
        \zihao { 3 } \bfseries \heiti
        \@@_line_spread:n { 1.5 }
      },
    bottom-skip = 1 cm,
  }
%    \end{macrocode}
% \end{macro}
%
% \begin{macro}[int, module=hust/component/]{decl/auth text}
% 学位论文版权使用授权书正文。
%    \begin{macrocode}
\@@_declare_component_instance:nnn { decl } { auth text }
  {
    content = { \c_@@_auth_decl_text_tl },
    format =
      {
        \zihao { -4 } \rmfamily
        \@@_line_spread:n { 1.5 }
      },
    align = justify,
    bottom-skip = 0.5 cm,
  }
%    \end{macrocode}
% \end{macro}
%
% \begin{macro}[int, module=hust/component/]{decl/auth secr}
% 学位论文保密使用保密选项。
%    \begin{macrocode}
\@@_declare_component_instance:nnn { decl } { auth secr }
  {
    content =
      {
        \mode_leave_vertical:
        \hbox_overlap_right:n
          { \clist_item:Nn \c_@@_auth_decl_secr_clist { 1 } }
        \skip_vertical:n { -1.2 cm }
        \skip_horizontal:n { 6 em }
        \clist_item:Nn \c_@@_auth_decl_secr_clist { 2 }
        \tex_par:D
        \skip_horizontal:n { 6 em }
        \clist_item:Nn \c_@@_auth_decl_secr_clist { 3 }
        \tex_par:D
        \clist_item:Nn \c_@@_auth_decl_secr_clist { 4 }
      },
    format =
      {
        \zihao { -4 } \rmfamily
        \@@_line_spread:n { 1.5 }
      },
    align = justify,
    bottom-skip = 1 cm,
  }
%    \end{macrocode}
% \end{macro}
%
% \begin{macro}[int, module=hust/component/]{decl/auth sign}
%    \begin{macrocode}
\@@_declare_component_instance:nnn { decl } { auth sign }
  {
    content =
      {
        \clist_map_inline:Nn \c_@@_auth_decl_sign_clist
          {
            #1
            \tex_par:D
          }
      },
    format =
      {
        \zihao { -4 } \rmfamily
        \@@_line_spread:n { 1.5 }
      },
    align = justify,
  }
%    \end{macrocode}
% \end{macro}
%
% \subsection{绘制封面}
%
% \begin{macro}{\makezhtitle}
% 中文标题页。
%    \begin{macrocode}
\@@_declare_page_instance:nn { title zh }
  {
    components = { identifiers, logo, level, degree, title, info, },
  }
\DeclareDocumentCommand \makezhtitle { }
  {
    \UseInstance { hust / page } { title zh }
  }
%    \end{macrocode}
% \end{macro}
%

% \begin{macro}{\makeentitle}
% 英文标题页。
%    \begin{macrocode}
\@@_declare_page_instance:nn { title en }
  {
    components = { degree, title, info, address, date, },
  }
\DeclareDocumentCommand \makeentitle { }
  {
    \UseInstance { hust / page } { title en }
  }
%    \end{macrocode}
% \end{macro}
%
% \begin{macro}{\makedecl}
% 声明页。
%    \begin{macrocode}
\@@_declare_page_instance:nn { decl }
  {
    components =
      {
        orig title, orig text, orig sign,
        auth title, auth text, auth secr, auth sign,
      },
  }
\DeclareDocumentCommand \makedecl { }
  {
    \UseInstance { hust / page } { decl }
  }
%    \end{macrocode}
% \end{macro}
%
% \begin{macro}{\maketitle}
% 封面。
%    \begin{macrocode}
\DeclareDocumentCommand \maketitle { }
  {
    \makezhtitle
    \makeentitle
    \makedecl
  }
%    \end{macrocode}
% \end{macro}
%
%
% \section{扉页样式}
%
% \subsection{答辩委员会页}
%
% \subsection{声明页}
%
% \subsection{目录样式}
%
% \section{正文样式}
%
% \subsection{章节样式}
%
%    \begin{macrocode}
\keys_set:nn { ctex }
  {
%    \end{macrocode}
% 设置章（chapter）、节（section）与小节（sub-section）标题样式。

%    \begin{macrocode}
    chapter =
      {
        format = \zihao{3} \bfseries \rmfamily \centering,
        name = {},
        beforeskip = 15.5 bp,
        afterskip = 14 bp,
        number = \arabic{ chapter },
        pagestyle = fancy,
      },
    section =
      {
        format = \zihao{4} \bfseries \rmfamily \raggedright,
        numberformat = \rmfamily,
        beforeskip = 10 bp,
        afterskip = 7 bp,
      },
    subsection =
      {
        format = \zihao{-4} \bfseries \rmfamily \raggedright,
        numberformat = \rmfamily,
        beforeskip = 7 bp,
        afterskip = 4 bp,
      }
  }
%    \end{macrocode}
%
% \section{篇后样式}
%
% \subsection{致谢页}
%
% \subsection{参考文献}
%
% \section{附录样式}
%
% \subsection{答辩委员会决议}
%
% \subsection{成果论文对应关系}
%
% \subsection{中英文缩写对照表}
%
% \subsection{程序代码样式}
%
% \section{文档类后处理}
%
% \section{用户接口}
%
% \begin{macro}[module=info/]{info/}
% \begin{macro}[module=style/]{style/}
% 定义元（meta）键值对。
%    \begin{macrocode}
\keys_define:nn { hust }
  {
    info .meta:nn = { hust / info } {#1},
    style .meta:nn = { hust / style } {#1}
  }
%    \end{macrocode}
% \end{macro}
% \end{macro}
%
% 文档类初始设置。学校代码默认为 \opt{10487}。
%    \begin{macrocode}
\keys_set:nn { hust }
  {
    info =
      {
        degree = { academic },
        secret = none,
        school-id = { 10487 },
        date =
          {
            \int_use:N \c_sys_year_int -
            \int_use:N \c_sys_month_int -
            \int_use:N \c_sys_day_int
          }
      }
  }
%    \end{macrocode}
%
% \begin{macro}{\hustsetup}
% 用户设置接口。
%    \begin{macrocode}
\NewDocumentCommand \hustsetup { m }
  { \keys_set:nn { hust } {#1} }
%    \end{macrocode}
% \end{macro}
%
%    \begin{macrocode}
%</class>
%    \end{macrocode}
%
% \section{常量定义}
%
%    \begin{macrocode}
%<*def-m|def-d>
%    \end{macrocode}
%
% \subsection{全局常量}
%
% \begin{variable}{\c_@@_name_school_zh_tl, \c_@@_name_school_en_tl}
% 学校名称。
%    \begin{macrocode}
\tl_const:Nn \c_@@_name_school_zh_tl { 华中科技大学 }
\tl_const:Nn \c_@@_name_school_en_tl { Huazhong~ University~ of~ Science~
  and~ Technology }
%    \end{macrocode}
% \end{variable}
%
% 常用标点符号，见表~\ref{tab:punctuations}，定义为 |\c_@@_ #1 _tl|。
%    \begin{macrocode}
\clist_map_inline:nn
  {
    { ideo_comma       } { ^^^^3001 },
    { ideo_full_stop   } { ^^^^3002 },
    { fwid_comma       } { ^^^^ff0c },
    { fwid_full_stop   } { ^^^^ff0e },
    { fwid_colon       } { ^^^^ff1a },
    { fwid_semicolon   } { ^^^^ff1b },
    { fwid_left_paren  } { ^^^^ff08 },
    { fwid_right_paren } { ^^^^ff09 }
  }
  { \@@_define_punct:nn #1 }
%    \end{macrocode}
%
% \begin{table}[ht]
%   \caption{常用标点符号}
%   \label{tab:punctuations}
%   \centering
%   \begin{tabular}{cccc}
%     \toprule
%       \textbf{中文名称} &
%       \textbf{英文名称} &
%       \textbf{符号} &
%       \textbf{Unicode 码位} \\
%     \midrule
%       中文顿号     & Ideographic comma           & ^^^^3001 & U+3001 \\
%       中文句号     & Ideographic full stop       & ^^^^3002 & U+3002 \\
%       中文逗号     & Fullwidth comma             & ^^^^ff0c & U+FF0C \\
%       全角西文句点 & Fullwidth full stop         & ^^^^ff0e & U+FF0E \\
%       中文冒号     & Fullwidth colon             & ^^^^ff1a & U+FF1A \\
%       中文分号     & Fullwidth semicolon         & ^^^^ff1b & U+FF1B \\
%       中文左圆括号 & Fullwidth left parenthesis  & ^^^^ff08 & U+FF08 \\
%       中文右圆括号 & Fullwidth right parenthesis & ^^^^ff09 & U+FF09 \\
%     \bottomrule
%   \end{tabular}
% \end{table}
%
% \begin{variable}{\c_@@_line_spread_fp}
% 设置行距倍数1.5倍。
%    \begin{macrocode}
\fp_const:Nn \c_@@_line_spread_fp { 1.5 }
%    \end{macrocode}
% \end{variable}
%
% \begin{variable}{\c_@@_emptybox_tl, \c_@@_checkbox_tl}
%    \begin{macrocode}
\tl_const:Nn \c_@@_emptybox_tl { \usym{2610} }
\tl_const:Nn \c_@@_checkbox_tl { \usym{2611} }
%    \end{macrocode}
% \end{variable}
%
% \subsection{封面页常量}
% \begin{variable}{
%   \c_@@_name_degree_font_tl,
%   \c_@@_name_professional_tl,
%   \c_@@_name_date_tl,
%   \c_@@_name_clc_tl,
%   \c_@@_name_student_id_tl,
%   \c_@@_name_school_id_tl,
%   \c_@@_name_secret_tl,
%   \c_@@_school_id_tl,
% }
% 中文封面页常量。
%    \begin{macrocode}
\clist_map_inline:nn
  {
%<def-m>    { degree_font } { \bfseries \songti },
%<def-d>    { degree_font } { \heiti },
    { date } { 答\hfill 辩\hfill 日\hfill 期 },
    { clc } { 分类号 },
    { student_id } { 学号 },
    { school_id } { 学校代码 },
    { secret } { 密级 },
  }
  { \@@_define_name:nn #1 }
\tl_const:Nn \c_@@_school_id_tl { 10487 }
%    \end{macrocode}
% \end{variable}
%
% \begin{variable}{
%   \c_@@_address_tl
% }
% 英文封面页常量。
%    \begin{macrocode}
\tl_const:Nn \c_@@_address_tl
  {
    Huazhong~ University~ of~ Science~ and~ Technology\\
    Wuhan~ 430074,~ P.~ R.~ China
  }
%    \end{macrocode}
% \end{variable}
%
% \begin{variable}{
%   \c_@@_name_level_en_tl,
%   \c_@@_name_level_zh_tl,
%   \c_@@_name_academic_en_tl,
%   \c_@@_name_academic_zh_tl,
%   \c_@@_name_professional_en_tl,
%   \c_@@_name_professional_zh_tl,
%   \c_@@_name_author_en_tl,
%   \c_@@_name_author_zh_tl,
%   \c_@@_name_major_en_tl,
%   \c_@@_name_major_zh_tl,
%   \c_@@_name_supervisor_en_tl,
%   \c_@@_name_supervisor_zh_tl,
% }
% 中英文页共用常量。
%    \begin{macrocode}
\clist_map_inline:nn
  {
    { level }
      { A~ Dissertation~ Submitted~ in~ Partial~ Fulfillment~ of~
%<def-m>      the~
%<def-d>      The~
      Requirements~ for~ }
%<def-m>      { 硕士学位论文 },
%<def-d>      { 博士学位论文 },
%<*def-m>
    { academic } { the~ Master~ Degree~ in } { 学术型 },
    { professional } { the~ Professional~ Master~ Degree~ in }  { 专业型 },
    { author } { Candidate } { 学\hfill 位\hfill 申\hfill 请\hfill 人 },
%</def-m>
%<*def-d>
    { academic } { The~ Degree~ of~ Doctor~ of } { 学术型 },
    { professional } { The~ Degree~ of~ Doctor~ of } { 专业型 }, % no diff in English
    { author } { Ph.D.~ Candidate } { 学\hfill 位\hfill 申\hfill 请\hfill 人 },
%</def-d>
    { major } { Major } { 学\hfill 科\hfill 专\hfill 业 },
    { supervisor } { Supervisor } { 指\hfill 导\hfill 老\hfill 师 },
  }
  { \@@_define_name:nnn #1 }
%    \end{macrocode}
% \end{variable}
%
% \subsection{声明页常量}
%
% \begin{variable}{\c_@@_orig_decl_title_tl, \c_@@_orig_decl_text_tl}
% 独创性声明。
%    \begin{macrocode}
\tl_const:Nn \c_@@_orig_decl_title_tl { 独创性声明 }
\tl_const:Nn \c_@@_orig_decl_text_tl
  {
    本人声明所呈交的学位论文是我个人在导师指导下进行的研究工作及取得的研
    究成果。尽我所知，除文中已经标明引用的内容外，本论文不包含任何其他个人或
    集体已经发表或撰写过的研究成果。对本文的研究做出贡献的个人和集体，均已在
    文中以明确方式标明。本人完全意识到本声明的法律结果由本人承担。
  }
%    \end{macrocode}
% \end{variable}
%
% \begin{variable}{\c_@@_orig_decl_sign_clist}
% 独创性声明签名项目。
%    \begin{macrocode}
\clist_const:Nn \c_@@_orig_decl_sign_clist
  {
    学位论文作者签名：\hspace{2em}\mbox{},
    日期：\hspace{2em}年\hspace{1.5em}月\hspace{1.5em}日
  }
%    \end{macrocode}
% \end{variable}
%
% \begin{variable}{\c_@@_auth_decl_title_tl, \c_@@_auth_decl_text_tl}
% 学位论文版权使用授权书。
%    \begin{macrocode}
\tl_const:Nn \c_@@_auth_decl_title_tl { 学位论文版权使用授权书 }
\tl_const:Nn \c_@@_auth_decl_text_tl
  {
    本学位论文作者完全了解学校有关保留、使用学位论文的规定，即：学校有权
    保留并向国家有关部门或机构送交论文的复印件和电子版，允许论文被查阅和借阅。
    本人授权华中科技大学可以将本学位论文的全部或部分内容编入有关数据库进行检
    索，可以采用影印、缩印或扫描等复制手段保存和汇编本学位论文。
  }
%    \end{macrocode}
% \end{variable}
%
% \begin{variable}{\c_@@_auth_decl_secr_clist, \c_@@_auth_decl_sign_clist}
% 学位论文版权使用授权书填写项目：保密与签名。
%    \begin{macrocode}
\clist_const:Nn \c_@@_auth_decl_secr_clist
  {
    本论文属于,
    保\hspace{1em}密□，在\underline{\hspace{3em}}年解密后适用本授权书。,
    不保密□。,
    （请在以上方框内打“√”）
  }
\clist_const:Nn \c_@@_auth_decl_sign_clist
  {
    学位论文作者签名：\hfill{} 指导教师签名：\hspace*{5em} \mbox{},
    日期：\hspace*{2em}年\hspace*{1.5em}月\hspace*{1.5em}日 \hfill{}
    日期：\hspace*{2em}年\hspace*{1.5em}月\hspace*{1.5em}日\hspace*{1em} \mbox{}
  }
%    \end{macrocode}
% \end{variable}
%
% \subsection{摘要页常量}
% \begin{variable}{
%   \c_@@_name_abstract_zh_tl,
%   \c_@@_name_abstract_en_tl,
%   \c_@@_name_keywords_zh_tl,
%   \c_@@_name_keywords_en_tl,
%   \c_@@_sep_keywords_zh_tl,
%   \c_@@_sep_keywords_en_tl,
% }
%    \begin{macrocode}
\clist_map_inline:nn
  {
    { abstract } { Abstract } { 摘要 },
    { keywords } { Key~ words: } { 关键词： },
  }
  { \@@_define_name:nnn #1 }
\tl_const:Nn \c_@@_sep_keywords_zh_tl { ； }
\tl_const:Nn \c_@@_sep_keywords_en_tl { , }
%    \end{macrocode}
% \end{variable}
%
% \subsection{正文常量}
% \begin{variable}{\c_@@_header_tl}
% 页眉。
%    \begin{macrocode}
%<def-m>\tl_const:Nn \c_@@_header_tl { 华中科技大学硕士学位论文 }
%<def-d>\tl_const:Nn \c_@@_header_tl { 华中科技大学博士学位论文 }
%    \end{macrocode}
% \end{variable}
%
%    \begin{macrocode}
%</def-m|def-d>
%    \end{macrocode}

% \section{文献样式}
%
% \subsection{标注样式}
%
%    \begin{macrocode}
%<*cbx>
%    \end{macrocode}
%
% 临时占位测试。
%    \begin{macrocode}
\RequireCitationStyle{gb7714-2015}
%    \end{macrocode}
%
%    \begin{macrocode}
%</cbx>
%    \end{macrocode}
%
% \subsection{著录样式}
%
%    \begin{macrocode}
%<*bbx>
%    \end{macrocode}
%
% 临时占位测试。
%    \begin{macrocode}
\RequireBibliographyStyle{gb7714-2015}
%    \end{macrocode}
%
%    \begin{macrocode}
%</bbx>
%    \end{macrocode}
%
% \section{文档样式}
%
% ^^A 本小节初版实现主要来自 \pkg{fduthesis} 文档类。
%
%    \begin{macrocode}
%<*doc-sty>
%<@@=hustdoc>
\RequirePackage
  {
    ulem,
    xcolor-material,
    zref-base,
  }
%    \end{macrocode}
%
% \subsection{内部函数}
%
% \begin{macro}{\@@_patch_cmd:Nnn, \@@_preto_cmd:Nn, \@@_appto_cmd:Nn}
% 补丁工具。
%    \begin{macrocode}
\cs_new_protected:Npn \@@_patch_cmd:Nnn #1#2#3
  {
    \ctex_patch_cmd_once:NnnnTF #1 { } {#2} {#3}
      { } { \ctex_patch_failure:N #1 }
  }
\cs_new_protected:Npn \@@_preto_cmd:Nn #1#2
  {
    \ctex_preto_cmd:NnnTF #1 { } {#2}
      { } { \ctex_patch_failure:N #1 }
  }
\cs_new_protected:Npn \@@_appto_cmd:Nn #1#2
  {
    \ctex_appto_cmd:NnnTF #1 { } {#2}
      { } { \ctex_patch_failure:N #1 }
  }
%    \end{macrocode}
% \end{macro}
%
% \subsection{\cls{l3doc} 补丁}
%
%    \begin{macrocode}
%<@@=codedoc>
%    \end{macrocode}
%
% \begin{macro}{\@@_typeset_functions:}
% 减少展开标记（用来表明中英文模板中的不同用法）前的空格。
%    \begin{macrocode}
\__hustdoc_patch_cmd:Nnn \@@_typeset_expandability:
  { & } { & \skip_horizontal:n { -0.5em } }
%    \end{macrocode}
% \end{macro}
%
% \begin{macro}{\@@_typeset_functions:, \@@_macro_init:,
%   \@@_macro_dump:}
% 左侧边注的函数列表采用单倍行距。
%    \begin{macrocode}
\__hustdoc_preto_cmd:Nn \@@_typeset_functions:
  { \MacroFont }
\__hustdoc_patch_cmd:Nnn \@@_macro_init:
  { \hbox:n } { \MacroFont \hbox:n }
\__hustdoc_patch_cmd:Nnn \@@_macro_dump:
  { \hbox_unpack_drop:N } { \MacroFont \hbox_unpack_drop:N }
%    \end{macrocode}
% \end{macro}
%
% \begin{macro}{\@@_macro_typeset_one:nN}
% 在 \env{macro} 环境的侧边栏中，\cls{l3doc} 根据命令的长短，分别用
% 普通字体和紧缩字体输出。然而很长的命令还是会超出页边。这里用缩放
% 盒子的手段使得长命令也可正常显示。
%    \begin{macrocode}
\cs_set_protected:Npn \@@_macro_typeset_one:nN #1#2
  {
    \vbox_set:Nn \l_@@_macro_box
      {
        \MacroFont
        \vbox_unpack_drop:N \l_@@_macro_box
        \hbox_set:Nn \l_tmpa_box
          { \@@_print_macroname:nN {#1} #2 }
%    \end{macrocode}
% \tn{marginparwidth} 和 \tn{marginparsep} 分别是边注的宽度及其到版心的距离，
% \tn{labelsep} 则是编号盒子右端与条目首行文本之间的距离。
%    \begin{macrocode}
        \dim_set:Nn \l_tmpa_dim
          { \marginparwidth - \labelsep - \marginparsep }
        \dim_compare:nNnT { \box_wd:N \l_tmpa_box } > \l_tmpa_dim
          {
            \box_resize_to_wd_and_ht:Nnn \l_tmpa_box
              { \l_tmpa_dim } { \box_ht:N \l_tmpa_box }
          }
        \hbox_overlap_left:n
          {
            \box_use:N \l_tmpa_box
            \skip_horizontal:n { \marginparsep - \labelsep }
          }
      }
    \int_incr:N \l_@@_macro_int
  }
%    \end{macrocode}
% \end{macro}
%
% \subsection{文档层命令}
%<@@=hustdoc>
%
% \begin{macro}[no-user-doc]{\key}
% 用于在文档中输出自定义索引的命令，主要用于键值对的键的索引。仅当索引项的父分类为 |option| 时，即用于文档类选项时，隐藏父分类。
%  \begin{arguments}
%    \item 在索引中显示项的父分类，默认为 \opt{option}。
%    \item 命令本体。
%  \end{arguments}
%    \begin{macrocode}
\DeclareDocumentCommand \key { O{option} m }
  {
    \tl_if_eq:nnTF { #1 } { option }
      { \cmd[module=#1/]{#2} }
      { \cmd[module=#1/]{#1/#2} }
  }
%    \end{macrocode}
% \end{macro}
%
% \begin{macro}[no-user-doc]{\opt}
% 输出键值对的值。
%    \begin{macrocode}
\DeclareDocumentCommand \opt { m }
  { \texttt{#1} }
%    \end{macrocode}
% \end{macro}
%
% \begin{macro}[int]{\TeX, \LaTeX, \LaTeXe,
%   \pdfTeX, \pdfLaTeX, \XeTeX, \XeLaTeX, \LuaTeX, \LuaLaTeX,
%   \AmSLaTeX, \TeXLive, \MiKTeX, \BibTeX, \biber, \TikZ
% }
% \TeX{} 相关标志。
%    \begin{macrocode}
\def\TeX{\hologo{TeX}}
\def\LaTeX{\hologo{LaTeX}}
\def\LaTeXe{\hologo{LaTeXe}}
\def\pdfTeX{\hologo{pdfTeX}}
\def\pdfLaTeX{\hologo{pdfLaTeX}}
\def\XeTeX{\hologo{XeTeX}}
\def\XeLaTeX{\hologo{XeLaTeX}}
\def\LuaTeX{\hologo{LuaTeX}}
\def\LuaLaTeX{\hologo{LuaLaTeX}}
\def\AmSLaTeX{\hologo{AmSLaTeX}}
\def\TeXLive{\TeX\ Live}
\def\MiKTeX{\hologo{MiKTeX}}
\def\BibTeX{\hologo{BibTeX}}
\def\biber{\hologo{biber}}
%    \end{macrocode}
% 该定义来自 \file{pgfmanual-en-macros.tex}。
%    \begin{macrocode}
\def\TikZ{Ti\emph{k}Z}
%    \end{macrocode}
% \end{macro}
%
% \subsection{\env{macrocode}环境}
%
% \subsubsection{代码行处理}
%
% \begin{macro}[int]{\xmacro@code,\sxmacro@code}
% 重新实现 \env{macrocode} 与 \env{macrocode*} 环境的核心功能，将对代码逐行
% 处理。后者会将空格显示为“\verb*| |”。
%    \begin{macrocode}
\cs_set_protected_nopar:Npn \xmacro@code
  { \@@_marco_code:w }
\cs_set_protected_nopar:Npn \sxmacro@code
  {
    \fontspec_print_visible_spaces:
    \xmacro@code
  }
%    \end{macrocode}
% \end{macro}
%
% \begin{macro}{\@@_marco_code:w}
%    \begin{macrocode}
\cs_new_protected_nopar:Npn \@@_marco_code:w
  {
%    \end{macrocode}
% 根据 \tn{codeline@index} 是否为 |true| 选择是否显示行号。
%    \begin{macrocode}
    \ifcodeline@index
      \@@_marco_code_every_par:n { \@@_code_line_no: }
    \else:
      \@@_marco_code_every_par:n { }
    \fi:
%    \end{macrocode}
% 设置代码段结束标记为“\verb*|%    \end{macrocode}^^M|”。
%    \begin{macrocode}
    \@@_make_finish_tag:x { \@currenvir }
%    \end{macrocode}
% 开始 \env{macrocode}。
%    \begin{macrocode}
    \@@_macro_code_start:w
  }
%    \end{macrocode}
% \end{macro}
%
% \begin{macro}{\@@_marco_code_every_par:n}
% 在每段之前插入内容。这里每段即相当于每行。
%    \begin{macrocode}
\cs_new_protected:Npn \@@_marco_code_every_par:n #1
  {
    \everypar
      {
        \everypar {#1}
        \if@inlabel
          \global \@inlabelfalse \@noparlistfalse
          \llap { \box \@labels \hskip \leftskip }
        \fi
        #1
      }
  }
%    \end{macrocode}
% \end{macro}
%
% 设置 \tn{endlinechar} 为 $-1$，表示行末不插入任何字符（实际上相当于在行尾
% 插入注释符 |%|）。
%    \begin{macrocode}
\group_begin:
  \int_set:Nn \tex_endlinechar:D { -1 }
%    \end{macrocode}
%
% \begin{variable}{\c_@@_active_space_tl}
% 活动字符类的空格（ASCII 码为 32）。
%    \begin{macrocode}
  \use:n
    {
      \char_set_catcode_active:n { 32 }
      \tl_const:Nn \c_@@_active_space_tl
    }
    { }
\group_end:
%    \end{macrocode}
% \end{variable}
%
% ASCII 码 13 是回车符 |^^M|。将其设置为活动字符。
%    \begin{macrocode}
\group_begin:
  \char_set_catcode_active:n { 13 }
%    \end{macrocode}
%
% \begin{macro}{\@@_make_finish_tag:n,\@@_make_finish_tag:x}
% \env{macrocode} 结尾标记。展开后变成“\verb*|%    \end{#1}^^M|”。
%    \begin{macrocode}
  \cs_new_protected:Npx \@@_make_finish_tag:n #1
    {
      \tl_set:Nn \exp_not:N \l_@@_macro_code_finish_tl
        {
          \c_percent_str
          \prg_replicate:nn { 4 }
            { \exp_not:o { \c_@@_active_space_tl } }
          \exp_not:o { \active@escape@char } end
          \c_left_brace_str #1 \c_right_brace_str
          \exp_not:N ^^M
        }
    }
  \cs_generate_variant:Nn \@@_make_finish_tag:n { x }
%    \end{macrocode}
% \end{macro}
%
% \begin{macro}{\@@_macro_code_start:w}
% 开始代码环境。此命令主要是为了防止 |\begin{macrocode}| 后出现多余的空行。
%    \begin{macrocode}
  \cs_new_protected:Npn \@@_macro_code_start:w #1
    {
      \str_if_eq:nnTF {#1} { ^^M }
        { \@@_macro_code_read_line:w }
        { \@@_macro_code_read_line:w #1 }
    }
%    \end{macrocode}
% \end{macro}
%
% \begin{macro}{\@@_macro_code_read_line:w}
% 逐行读取代码，并连同行尾回车符一并存入 \cs{l_@@_macro_code_line_tl}。如果该行
% 与结束标记“\verb*|%    \end{macrocode}^^M|”相同，则结束此 \env{macrocode}；
% 否则继续处理该行代码。
%    \begin{macrocode}
  \cs_new_protected:Npn \@@_macro_code_read_line:w #1 ^^M
    {
      \tl_set:Nn \l_@@_macro_code_line_tl { #1 ^^M }
      \tl_if_eq:NNTF
        \l_@@_macro_code_line_tl \l_@@_macro_code_finish_tl
        { \exp_args:Nx \end { \@currenvir } }
        {
          \@@_macro_code_process_line:
          \@@_macro_code_read_line:w
        }
    }
%    \end{macrocode}
% \end{macro}
%
% \begin{macro}{\@@_swap_cr:,\@@_swap_cr_aux:w}
% 把 \cs{l_@@_macro_code_line_tl} 中的回车符 |^^M| 挪到外面。
%    \begin{macrocode}
  \cs_new_protected:Npn \@@_swap_cr:
    {
      \exp_after:wN
        \@@_swap_cr_aux:w \l_@@_macro_code_line_tl
    }
  \cs_new_protected:Npn \@@_swap_cr_aux:w #1 ^^M
    {
      \group_insert_after:N ^^M
      \tl_set:Nn \l_@@_macro_code_line_tl {#1}
    }
%    \end{macrocode}
% \end{macro}
%
% \begin{variable}{\c_@@_active_cr_tl}
% 活动字符类的回车符。
%    \begin{macrocode}
  \tl_const:Nn \c_@@_active_cr_tl { ^^M }
\group_end:
%    \end{macrocode}
% \end{variable}
%
% \begin{variable}{\l_@@_macro_code_line_tl,
%   \l_@@_macro_code_finish_tl,
%   \g_@@_macro_code_verbatim_stop_tl}
% 分别用来存储代码行、\env{macrocode} 结束标记以及抄录停止标记。
%    \begin{macrocode}
\tl_new:N \l_@@_macro_code_line_tl
\tl_new:N \l_@@_macro_code_finish_tl
\tl_new:N \g_@@_macro_code_verbatim_stop_tl
%    \end{macrocode}
% \end{variable}
%
% \begin{macro}{\@@_process_normal_line:}
% 普通代码行根据开头字符分别处理。
%    \begin{macrocode}
\cs_new_protected_nopar:Npn \@@_process_normal_line:
  {
    \str_case_e:nnF
      { \str_head:N \l_@@_macro_code_line_tl }
      {
%    \end{macrocode}
% 以 |%| 开头的行先由 \cs{tl_tail:N} 去掉 |%|，之后再检查 |<|。
%    \begin{macrocode}
        { \c_percent_str }
        {
          \@@_check_angle:x
            { \tl_tail:N \l_@@_macro_code_line_tl }
        }
%    \end{macrocode}
% 以 |#| 开头的行按注释的格式输出。
%    \begin{macrocode}
        { \c_hash_str }
        { \@@_output_comment_line: }
      }
%    \end{macrocode}
% 其余正常输出。
%    \begin{macrocode}
      { \@@_output_line: }
  }
%    \end{macrocode}
% \end{macro}
%
% \begin{macro}{\@@_process_verbatim_line:}
% 处理抄录代码行（不检查 |%| 与 |<|）。
%    \begin{macrocode}
\cs_new_protected_nopar:Npn \@@_process_verbatim_line:
  {
%    \end{macrocode}
% 将该行与抄录停止标记进行比较。
%    \begin{macrocode}
    \tl_if_eq:NNTF \l_@@_macro_code_line_tl
        \g_@@_macro_code_verbatim_stop_tl
%    \end{macrocode}
% 若相同，则结束抄录环境，清空抄录停止标记，并输出该标记；
%    \begin{macrocode}
      {
        \tl_gclear:N \g_@@_macro_code_verbatim_stop_tl
        \cs_gset_eq:NN \@@_macro_code_process_line:
          \@@_process_normal_line:
        \@@_output_module:nn
          { \color { verb@guard } }
          {
            \@@_swap_cr:
            \@@_module_pop:n { \l_@@_macro_code_line_tl }
          }
      }
%    \end{macrocode}
% 否则直接输出抄录代码。
%    \begin{macrocode}
      { \tl_use:N \l_@@_macro_code_line_tl }
  }
%    \end{macrocode}
% \end{macro}
%
% \begin{macro}{\@@_macro_code_process_line:}
% 处理代码行。该命令的作用如下：
% \begin{itemize}
%   \item 正常情况下，等同于 \cs{@@_process_normal_line:}；
%   \item 在 \cs{DontCheckModules} 之后，等价于 \cs{@@_output_line:}，
%     即不检查模块标记，直接输出；
%   \item 在抄录环境中，等价于 \cs{@@_process_verbatim_line:}，
%     此时将不再特殊处理以 |%| 开头的代码行。
% \end{itemize}
%    \begin{macrocode}
\cs_new_eq:NN \@@_macro_code_process_line:
  \@@_process_normal_line:
%    \end{macrocode}
% \end{macro}
%
% \subsubsection{模块标记处理}
%
% \begin{macro}[int]{\CheckModules, \DontCheckModules}
% 选择是否检查模块标记。这两个命令在 \pkg{doc} 宏包中已有定义，此处重新声明。
%    \begin{macrocode}
\DeclareDocumentCommand \CheckModules { }
  {
    \cs_set_eq:NN \@@_macro_code_process_line:
      \@@_process_normal_line:
  }
\DeclareDocumentCommand \DontCheckModules { }
  {
    \cs_set_eq:NN \@@_macro_code_process_line:
      \@@_output_line:
  }
%    \end{macrocode}
% \end{macro}
%
% \begin{macro}{\@@_check_angle:n,\@@_check_angle:x}
% 检查第一个字符是否是 |<|。若是，则检查模块；否则立刻输出改行。
% 该函数的参数不带 |%|。
%    \begin{macrocode}
\cs_new_protected:Npn \@@_check_angle:n #1
  {
    \str_if_eq:eeTF { \str_head:n {#1} } { < }
      { \@@_check_module:x { \tl_tail:n {#1} } }
      { \@@_output_comment_line: }
  }
\cs_generate_variant:Nn \@@_check_angle:n { x }
%    \end{macrocode}
% \end{macro}
%
% \begin{macro}{\@@_check_module:n,\@@_check_module:x}
% 检查紧跟 |<| 之后的字符。共有四种情况：
% \begin{itemize}
%   \item |*|：模块开始（|%<*|\meta{module}|>|）；
%   \item |/|：模块结束（|%</|\meta{module}|>|）；
%   \item |@|：名字空间（|%<@@=|\meta{namespace}|>|）；
%   \item |<|：抄录环境开始（|%<<|\meta{end-tag}）。
% \end{itemize}
% 若不是这几种情况，则为单独一行的独立模块（|%<|\meta{module}|>|）。
%    \begin{macrocode}
\cs_new_protected:Npn \@@_check_module:n #1
  {
    \str_case_e:nnF { \str_head:n {#1} }
      {
        { * } { \@@_module_star:w }
        { / } { \@@_module_slash:w }
        { @ } { \@@_module_at:w }
        { < } { \@@_module_verb:w }
      }
      { \@@_module_pm:w }
%    \end{macrocode}
% 参数 |#1| 将被上面几个 |:w| 型函数吃掉。
%    \begin{macrocode}
    #1 \q_stop
  }
\cs_generate_variant:Nn \@@_check_module:n { x }
%    \end{macrocode}
% \end{macro}
%
% 设置 |>| 为活动字符。
%    \begin{macrocode}
\group_begin:
  \char_set_catcode_active:N \>
%    \end{macrocode}
%
% \begin{macro}{\@@_module_star:w}
% 模块开始标记。
% \begin{arguments}
%   \item |*|\meta{module}
%   \item 之后的代码
% \end{arguments}
%    \begin{macrocode}
  \cs_new_protected:Npn \@@_module_star:w #1 > #2 \q_stop
    {
%    \end{macrocode}
% 临时变量 \cs{l_@@_tmpa_tl} 保存 |<*|\meta{module}|>| 之后的部分，即真实代码。
%    \begin{macrocode}
      \tl_set:Nn \l_@@_tmpa_tl {#2}
%    \end{macrocode}
% 判断该行是否为空（只含一个回车符 |^^M|）。
%    \begin{macrocode}
      \tl_if_eq:NNTF \l_@@_tmpa_tl \c_@@_active_cr_tl
%    \end{macrocode}
% 若是，则将 |<|\meta{module}|>| 放在行号的右侧；
%    \begin{macrocode}
        {
          \@@_output_module:nn
            { \@@_star_color: }
            {
              \@@_module_push:n
                { \@@_module_angle:n {#1} }
            }
        }
%    \end{macrocode}
% 否则放在左侧，并输出相应代码。
%    \begin{macrocode}
        {
          \@@_output_module_left:nn
            { \@@_star_color: }
            {
              \@@_module_push:n
                { \@@_module_angle:n {#1} }
            }
        }
      \@@_output_line:n {#2}
    }
%    \end{macrocode}
% \end{macro}
%
% \begin{macro}{\@@_module_slash:w}
% 模块结束标记。结构与 \cs{@@_module_star:w} 相同。
% \begin{arguments}
%   \item |/|\meta{module}
%   \item 之后的代码
% \end{arguments}
%    \begin{macrocode}
  \cs_new_protected:Npn \@@_module_slash:w #1 > #2 \q_stop
    {
      \tl_set:Nn \l_@@_tmpa_tl {#2}
      \tl_if_eq:NNTF \l_@@_tmpa_tl \c_@@_active_cr_tl
        {
          \@@_output_module:nn
            { \@@_slash_color: }
            {
              \@@_module_pop:n
                { \@@_module_angle:n {#1} }
            }
        }
        {
          \@@_output_module_left:nn
            { \@@_slash_color: }
            {
              \@@_module_pop:n
                { \@@_module_angle:n {#1} }
            }
        }
      \@@_output_line:n {#2}
    }
%    \end{macrocode}
% \end{macro}
%
% \begin{macro}{\@@_module_at:w}
% 名字空间。
% \begin{arguments}
%   \item 名字空间的名称（不含 |@@=|）
%   \item 之后的代码
% \end{arguments}
%    \begin{macrocode}
  \cs_new_protected:Npn \@@_module_at:w @ @ = #1 > #2 \q_stop
    {
      \@@_output_module:nn
        { \color { at@guard } }
        { \@@_module_angle:n { @ @ = #1 } }
%    \end{macrocode}
% 设置名字空间为 |#1|。\cls{l3doc} 中将名字空间称为“模块”（module），
% 注意不要混淆。
%    \begin{macrocode}
      \tl_gset:Nn \g__codedoc_module_name_tl {#1}
      \@@_output_line:n {#2}
    }
%    \end{macrocode}
% \end{macro}
%
% \begin{macro}{\@@_module_verb:w}
% 抄录开始。|#1|: |<|\meta{end-tag}，只有一个 |<|。\meta{end-tag} 的最后会
% 带有一个回车符 |^^M|。
%    \begin{macrocode}
  \cs_new_protected:Npn \@@_module_verb:w #1 \q_stop
    {
%    \end{macrocode}
% 重定义 \cs{@@_macro_code_process_line:}，用以输出抄录行。
%    \begin{macrocode}
      \cs_gset_eq:NN \@@_macro_code_process_line:
        \@@_process_verbatim_line:
%    \end{macrocode}
% 设置抄录停止标记。用 \cs{tl_tail:n} 去掉开头的 |<|。
%    \begin{macrocode}
      \tl_gset:Nx \g_@@_macro_code_verbatim_stop_tl
        { \c_percent_str \tl_tail:n {#1} }
%    \end{macrocode}
% 输出 |%<<|\meta{end-tag}。
%    \begin{macrocode}
      \@@_output_module:nn
        { \color { verb@guard } }
        {
          \@@_swap_cr:
          \@@_module_push:n { \l_@@_macro_code_line_tl }
        }
    }
%    \end{macrocode}
% \end{macro}
%
% \begin{macro}{\@@_module_pm:w}
% 处理单独一行的模块。|<|\meta{module}|>| 放在行号的左侧。
% \begin{arguments}
%   \item \meta{module}
%   \item 之后的代码
% \end{arguments}
%    \begin{macrocode}
  \cs_new_protected:Npn \@@_module_pm:w #1 > #2 \q_stop
    {
      \@@_output_module_left:nn
        { \@@_pm_color: }
        { \@@_module_angle:n {#1} }
      \@@_output_line:n {#2}
    }
\group_end:
%    \end{macrocode}
% \end{macro}
%
% \begin{macro}{\@@_output_line:n,\@@_output_line:}
% 输出代码行。参数将被存入 \cs{l_@@_macro_code_line_tl}，
% 再由不带参数的版本调用。
%    \begin{macrocode}
\cs_new_protected:Npn \@@_output_line:n #1
  {
    \tl_set:Nn \l_@@_macro_code_line_tl {#1}
%    \end{macrocode}
% 若为空行（只含一个 |^^M|），则直接输出（换行）。
%    \begin{macrocode}
    \tl_if_eq:NNTF
      \l_@@_macro_code_line_tl \c_@@_active_cr_tl
      { \tl_use:N \l_@@_macro_code_line_tl }
      {
%    \end{macrocode}
% 检查开头是否为 |%|，据此分别处理。
%    \begin{macrocode}
        \str_if_eq:eeTF
          { \str_head:N \l_@@_macro_code_line_tl } { \c_percent_str }
          { \@@_output_comment_line: } { \@@_output_line: }
      }
  }
\cs_new_protected_nopar:Npn \@@_output_line:
  {
    \tex_noindent:D
%    \end{macrocode}
% 此处将把 |@@| 替换为相应的名字空间。
%    \begin{macrocode}
    \@@_replace_at_at:N \l_@@_macro_code_line_tl
    \tl_use:N \l_@@_macro_code_line_tl
  }
%    \end{macrocode}
% \end{macro}
%
% \begin{macro}{\@@_output_comment_line:}
% 输出注释代码行。用灰色、斜体显示。
%    \begin{macrocode}
\cs_new_protected:Npn \@@_output_comment_line:
  {
    \tex_noindent:D
    \group_begin:
      \__hustdoc_swap_cr:
      \color { code@gray } \slshape \@@_output_line:
    \group_end:
  }
%    \end{macrocode}
% \end{macro}
%
% \begin{macro}{\@@_replace_at_at:N,
%   \@@_replace_at_at_aux:Nn,\@@_replace_at_at_aux:No}
% 把 |@@| 替换为相应的名字空间。其名称存放在全局变量
% \cs{g__codedoc_module_name_tl} 中。它为空时（|%<@@=>|），不做替换。
%    \begin{macrocode}
\cs_new_protected:Npn \@@_replace_at_at:N #1
  {
    \tl_if_empty:NF \g__codedoc_module_name_tl
      { \@@_replace_at_at_aux:No #1 \g__codedoc_module_name_tl }
  }
\cs_new_protected:Npn \@@_replace_at_at_aux:Nn #1#2
  {
%    \end{macrocode}
% 以下代码分别对应两种名字空间的替换：
% \begin{itemize}
%   \item 内部变量：|\|\meta{type}|_@@_|\meta{name} $\to$
%     |\|\meta{type}|__|\meta{namespace}|_|\meta{name}；
%   \item 内部函数：|\@@_|\meta{name}” $\to$
%     |\__|\meta{namespace}|_|\meta{name}”）。
% \end{itemize}
%    \begin{macrocode}
    \tl_replace_all:Nnn #1 { _ @ @ } { _ _ #2 }
    \tl_replace_all:Nnn #1 { @ @ } { _ _ #2 }
  }
\cs_generate_variant:Nn \@@_replace_at_at_aux:Nn { No }
%    \end{macrocode}
% \end{macro}
%
% \begin{macro}{\@@_module_push:n,
%   \@@_module_push_aux:nn,\@@_module_push_aux:on}
% 将模块名压入栈中。此处的栈主要用来处理模块名（包括抄录标记）之间的超链接。
%    \begin{macrocode}
\cs_new_protected_nopar:Npn \@@_module_push:n
  { \@@_module_push_aux:on { \int_use:N \c@HD@hypercount } }
\cs_new_protected:Npn \@@_module_push_aux:nn #1
  {
    \seq_gpush:Nn \g_@@_module_dest_seq {#1}
    \hypersetup { hidelinks }
    \exp_args:Nx \hdclindex
      { \zref@extractdefault { HD.#1 } { guard@end } { 1 } } { }
  }
\cs_generate_variant:Nn \@@_module_push_aux:nn { on }
%    \end{macrocode}
% \end{macro}
%
% \begin{macro}{\@@_module_pop:n,
%   \@@_module_pop_aux:nn,\@@_module_pop_aux:on}
% 将模块名释放出栈。
%    \begin{macrocode}
\cs_new_protected_nopar:Npn \@@_module_pop:n
  {
    \seq_gpop:NNTF \g_@@_module_dest_seq \l_@@_tmpa_tl
      { \@@_module_pop_aux:on { \l_@@_tmpa_tl } }
      { \BOOM \use:n }
  }
\cs_new_protected:Npn \@@_module_pop_aux:nn #1
  {
    \zref@labelbylist { HD.#1 } { hustdoc }
    \hypersetup { hidelinks }
    \hdclindex {#1} { }
  }
\cs_generate_variant:Nn \@@_module_pop_aux:nn { on }
%    \end{macrocode}
% \end{macro}
%
% \begin{variable}{\g_@@_module_dest_seq}
% 存放模块名的序列（栈）。
%    \begin{macrocode}
\seq_new:N \g_@@_module_dest_seq
%    \end{macrocode}
% \end{variable}
%
% 处理行号超链接。使用 \pkg{zref} 宏包。
%    \begin{macrocode}
\zref@newlist { hustdoc }
\zref@newprop { guard@end } [ 1 ]
  { \int_eval:n { \c@HD@hypercount - 1 } }
\zref@addprop { hustdoc } { guard@end }
%    \end{macrocode}
%
% \subsubsection{格式处理}
%
% \begin{macro}[int]{\MacroFont}
% 代码部分的字体。
%    \begin{macrocode}
\cs_set_protected:Npn \MacroFont
  {
    \linespread { 1.05 }
    \small \ttfamily \mdseries \upshape
  }
%    \end{macrocode}
% \end{macro}
%
% \begin{macro}{\@@_output_module:nn,\@@_output_module_left:nn}
% 输出模块名（分行内和行号左侧两种）。
% \begin{arguments}
%   \item 颜色等样式
%   \item 模块名
% \end{arguments}
%    \begin{macrocode}
\cs_new_protected:Npn \@@_output_module:nn #1#2
  {
    \tex_noindent:D
    \group_begin:
      #1
      \footnotesize \normalfont \sffamily #2
    \group_end:
  }
\cs_new_protected:Npn \@@_output_module_left:nn #1#2
  {
    \tex_noindent:D
    \hbox_overlap_left:n
      {
        \@@_output_module:nn {#1} {#2}
        \skip_horizontal:n { \leftskip + \smallskipamount }
      }
  }
%    \end{macrocode}
% \end{macro}
%
% \begin{macro}{\@@_star_color:,\@@_slash_color:,\@@_pm_color:}
% 选择模块标记的颜色。模块标记的颜色会根据嵌套层次而改变。
%    \begin{macrocode}
\cs_new_protected_nopar:Npn \@@_star_color:
  {
    \seq_gpop:NNTF \g_@@_star_color_seq \current@color
      { \set@color }
      { \@@_select_color: }
    \seq_gpush:No \g_@@_slash_color_seq { \current@color }
  }
\cs_new_protected_nopar:Npn \@@_slash_color:
  {
    \seq_gpop:NNTF \g_@@_slash_color_seq \current@color
      {
        \set@color
        \seq_gpush:No \g_@@_star_color_seq { \current@color }
      }
% TODO: 需要报错：star 与 slash 没有匹配
      { \BOOM }
  }
\cs_new_protected_nopar:Npn \@@_pm_color:
  {
    \seq_get:NNTF \g_@@_star_color_seq \current@color
      { \set@color }
      {
        \@@_select_color:
        \seq_gpush:No \g_@@_star_color_seq { \current@color }
      }
  }
%    \end{macrocode}
% \end{macro}
%
% \begin{variable}{\g_@@_star_color_seq,\g_@@_slash_color_seq}
% 存放模块标记颜色的序列。
%    \begin{macrocode}
\seq_new:N \g_@@_star_color_seq
\seq_new:N \g_@@_slash_color_seq
%    \end{macrocode}
% \end{variable}
%
% \begin{macro}{\@@_select_color:}
% \begin{macro}[int]{guard@series}
% 设置模块标记的色系。
%    \begin{macrocode}
\cs_new_protected_nopar:Npn \@@_select_color:
  { \color { guard@series!!+ } }
\definecolorseries { guard@series }
  { cmyk } { last } { blue } { purple }
%    \end{macrocode}
% \end{macro}
% \end{macro}
%
% 设置色系的增量大小。可选参数 |3| 意味着基色（blue）与末色（purple）之间将被
% 等分为三份。该数字比嵌套最大深度小 1。
%    \begin{macrocode}
\resetcolorseries [ 3 ] { guard@series }
%    \end{macrocode}
%
% \begin{macro}[int]{verb@guard,at@guard,code@gray}
% 设置颜色。
%    \begin{macrocode}
\definecolor { verb@guard } { named } { MaterialLime600 }
\definecolor { at@guard } { named } { MaterialPink }
\definecolor { code@gray } { named } { MaterialGrey }
%    \end{macrocode}
% \end{macro}
%
% \begin{macro}{\@@_module_angle:n}
% 输出“$\langle\cdots\rangle$”。
%    \begin{macrocode}
\cs_new_protected:Npn \@@_module_angle:n #1
  { < #1 > }
%   { \textlangle #1 \textrangle }
%   { \ensuremath \langle #1 \ensuremath \rangle }
%    \end{macrocode}
% \end{macro}
%
% \begin{macro}{\@@_code_line_no:}
% 行号。设置为阿拉伯数字。
%    \begin{macrocode}
\cs_new_protected_nopar:Npn \@@_code_line_no:
  {
    \int_gincr:N \c@CodelineNo
    \hbox_overlap_left:n
      {
        \hbox_to_wd:nn
          { \MacroIndent }
          {
            \HD@target
            \tex_hss:D \@@_code_line_no_style:
            \theCodelineNo \enspace
          }
        \tex_kern:D \@totalleftmargin
      }
  }
\tl_set:Nn \theCodelineNo { \arabic { CodelineNo } }
%    \end{macrocode}
% \end{macro}
%
% \begin{macro}{\@@_code_line_no_style:}
% 行号格式。
%    \begin{macrocode}
\cs_new_protected_nopar:Npn \@@_code_line_no_style:
  { \color { code@gray } \normalfont \sffamily \tiny }
%    \end{macrocode}
% \end{macro}
%
% \begin{macro}[int]{\HD@SetMacroIndent}
% 设置代码缩进（行号一栏的宽度）。该命令会写进 |.aux| 辅助文件，
% 以便在二次编译时确定最大行号宽度。
%    \begin{macrocode}
\cs_set_protected:Npn \HD@SetMacroIndent #1
  {
    \group_begin:
      \settowidth \MacroIndent
        {
          \@@_code_line_no_style:
          \prg_replicate:nn { \tl_count:n {#1} } { 0 }
          \enspace
        }
      \dim_gset_eq:NN \MacroIndent \MacroIndent
    \group_end:
  }
%    \end{macrocode}
% \end{macro}
%

% \subsection{修订记录}
%
% \begin{macro}{\@@_ltx_changes:nnn}
% 保存 \pkg{doc} 中 \tn{changes@} 的定义。
%    \begin{macrocode}
\cs_new_eq:NN \@@_ltx_changes:nnn \changes@
%    \end{macrocode}
% \end{macro}
%
% \begin{macro}[int]{\changes@}
% 重定义 \tn{changes@}，在版本号一行显示修改日期。
%    \begin{macrocode}
\cs_set_protected:Npn \changes@ #1#2
  {
    \@@_save_version_date:nn {#1} {#2}
    \@@_ltx_changes:nnn {#1} {#2}
  }
%    \end{macrocode}
% \end{macro}
%
% \begin{variable}{\g_@@_version_date_prop}
% 存放版本号与对应的修改日期。其中键为版本号，值为 |{开始日期，结束日期}|。
% 开始日期与结束日期可以相同。
%    \begin{macrocode}
\prop_new:N \g_@@_version_date_prop
%    \end{macrocode}
% \end{variable}
%
% \begin{macro}{\@@_save_version_date:nn}
% |nn| 版本最终将被 \tn{changes@} 调用。
% \begin{arguments}
%   \item 版本号
%   \item 日期
% \end{arguments}
% 它们分别对应 \tn{change} 的前两个参数（第三个是说明文字）。
%    \begin{macrocode}
\cs_new_protected:Npn \@@_save_version_date:nn #1#2
  {
    \prop_get:NnNTF \g_@@_version_date_prop
      {#1} \l_@@_tmpa_tl
      {
%    \end{macrocode}
% \cs{l_@@_tmpa_tl} 相当于两个参数（开始日期、结束日期），因此需要提前展开。
%    \begin{macrocode}
        \exp_after:wN
          \@@_save_version_date_aux:nnnn \l_@@_tmpa_tl
        {#2} {#1}
      }
      { \@@_save_version_date_aux:nnn {#1} {#2} {#2} }
  }
%    \end{macrocode}
% \end{macro}
%
% \begin{macro}{\@@_save_version_date_aux:nnnn}
% \begin{arguments}
%   \item 原开始日期
%   \item 原结束日期（显然应有 |#1| < |#2|）
%   \item 新读入的日期
%   \item 版本号
% \end{arguments}
% 如果 |#3| < |#1|，则读入日期 |#3|、|#2|；
% 如果 |#3| > |#2|，则读入日期 |#1|、|#3|。
%    \begin{macrocode}
\cs_new_protected:Npn \@@_save_version_date_aux:nnnn #1#2#3#4
  {
    \@@_if_date_later:nnTF {#1} {#3}
      { \@@_save_version_date_aux:nnn {#4} {#3} {#2} }
      {
        \@@_if_date_later:nnT {#3} {#2}
          { \@@_save_version_date_aux:nnn {#4} {#1} {#3} }
      }
  }
%    \end{macrocode}
% \end{macro}
%
% \begin{macro}{\@@_save_version_date_aux:nnn}
% 将版本号和日期存入 \cs{g_@@_version_date_prop}。
% \begin{arguments}
%   \item 版本号
%   \item 开始日期
%   \item 结束日期
% \end{arguments}
%    \begin{macrocode}
\cs_new_protected:Npn \@@_save_version_date_aux:nnn #1#2#3
  { \prop_gput:Nnn \g_@@_version_date_prop {#1} { {#2} {#3} } }
%    \end{macrocode}
% \end{macro}
%
% \begin{macro}[TF]{\@@_if_date_later:nn}
% \begin{macro}{\@@_parse_date:w}
% 比较两个日期。如果 |#1| 在 |#2| 之后，则为 |true|；反之为 |false|。
% 日期的格式为 |YYYY-MM-DD|。比较方法是直接将日期化成 8 位数字，
% 所以月、日前的 0 不可以省略。
%    \begin{macrocode}
\prg_new_conditional:Npnn \@@_if_date_later:nn #1#2 { TF, T }
  {
    \if_int_compare:w
        \@@_parse_date:w #1 / / / 0 \q_stop >
        \@@_parse_date:w #2 / / / 0 \q_stop \exp_stop_f:
      \prg_return_true:
    \else:
      \prg_return_false:
    \fi:
  }
\cs_new:Npn \@@_parse_date:w #1-#2-#3/ #4 \q_stop
  { #1#2#3 }
%    \end{macrocode}
% \end{macro}
% \end{macro}
%
% \begin{macro}[int]{\CTEX@versionitem}
% 版本条目标签。如果版本号不在 \cs{g_@@_version_date_prop} 里面，
% 则利用未定义的 \cs{BOOM} 报错。
%    \begin{macrocode}
\cs_new_protected:Npn \CTEX@versionitem #1 \efill
  {
    \@idxitem
    \prop_get:NnNTF \g_@@_version_date_prop
      {#1} \l_@@_tmpa_tl
      {
        \exp_after:wN
          \@@_print_version_date:nnn \l_@@_tmpa_tl
        {#1}
      }
      { \BOOM }
  }
%    \end{macrocode}
% \end{macro}
%
% \begin{macro}{\@@_print_version_date:nnn}
% 输出版本号和日期。如果开始日期和结束日期相同，则只输出一项。
% \begin{arguments}
%   \item 开始日期
%   \item 结束日期
%   \item 版本号
% \end{arguments}
%    \begin{macrocode}
\cs_new_protected:Npn \@@_print_version_date:nnn #1#2#3
  {
    \noindent
    \Hy@raisedlink { \belowpdfbookmark {#3} { HD.#3 } }
    \textbf {#3} \hfill
    \hbox:n
      {
        \footnotesize
        \str_if_eq:nnTF {#1} {#2}
          { ( #1 ) } { ( #1 \c_space_tl \c_tilde_str \c_space_tl #2 ) }
      }
    \par \nopagebreak
  }
%    \end{macrocode}
% \end{macro}
%
% \begin{macro}[int]{\HDorg@theglossary}
% 该命令由 \pkg{hypdoc} 宏包定义，用于存放标准文档类 \cls{book} 中定义的
% \tn{theindex} 命令。此处的补丁将在版本号一行最后加上修改日期。
%    \begin{macrocode}
\ctex_patch_cmd:Nnn \HDorg@theglossary
  { \let \item \@idxitem }
  { \let \item \CTEX@versionitem }
%    \end{macrocode}
% \end{macro}
%
% \begin{macro}[int]{\@wrglossary}
% 该命令由 \LaTeXe{} 内核定义，又由 \pkg{hypdoc} 宏包作了修改。
% 此处的补丁使得修订记录条目的页码能够指向对应行。
%    \begin{macrocode}
\ctex_patch_cmd:Nnn \@wrglossary
  { hdpindex }
  {
    \ifnum \c@HD@hypercount = \z@
      hdpindex
    \else
      hdclindex { \the \c@HD@hypercount }
    \fi
  }
%    \end{macrocode}
% \end{macro}
%
%    \begin{macrocode}
%</doc-sty>
%    \end{macrocode}
%
% \end{implementation}
